\documentclass[11pt]{article}
\usepackage[margin=1.1in]{geometry}
\usepackage{amsmath,amssymb,amsthm,mathtools}
\usepackage{booktabs}
\usepackage{xcolor}
\usepackage[colorlinks=true,linkcolor=blue!50!black,citecolor=blue!50!black,urlcolor=blue!50!black]{hyperref}
\usepackage{enumitem}
\newtheorem{theorem}{Theorem}[section]
\newtheorem{proposition}[theorem]{Proposition}
\newtheorem{lemma}[theorem]{Lemma}
\newtheorem{corollary}[theorem]{Corollary}
\newtheorem{conjecture}[theorem]{Conjecture}
\newtheorem{hypothesis}[theorem]{Hypothesis}
\theoremstyle{definition}
\newtheorem{definition}[theorem]{Definition}
\newtheorem{remark}[theorem]{Remark}
\DeclareMathOperator{\sgn}{sgn}
\newcommand{\dtv}{d_{\mathrm{TV}}}
\newcommand{\E}{\mathbb{E}}
\renewcommand{\Pr}{\mathbb{P}}
\newcommand{\DD}{\mathcal{D}}
\newcommand{\PP}{P}
\newcommand{\QQ}{Q}
\newcommand{\CC}{C}

\title{The cycle type of two rows of a random Latin square:\\ exact identities and a reduction of the weak Cavenagh--Greenhill--Wanless conjecture}
\author{[author]}
\date{3 October 2026 --- draft}

\begin{document}
\maketitle

\begin{abstract}
Let $L$ be a uniformly random Latin square of order $n$ and $\sigma$ the permutation of columns
carrying its first row to its second.  Cavenagh, Greenhill and Wanless conjectured that the
cycle type of $\sigma$ is asymptotically that of a uniform derangement, in total variation.
We prove three exact identities and a reduction theorem.  The marked-pair identity expresses the ratio of
completion counts of adjacent cycle types as the odds of one bit (``rows $1,2$ lie in a
common column cycle of a marked column pair'') in a single uniform square; a reduction
theorem, by which the conjecture follows from that bit being fair to $o(1/\log n)$; and
the ladder identity, by which the bit is exactly fair except on the squares in which no
``ladder pair'' of rows is flippable.  An orbit method bounds the probability of
the exceptional set by a product of conditionally independent factors.  The
conjecture is thereby reduced to a ``witness lemma'' (Hypothesis~\ref{hyp:witness}): that
with probability $1-o(1/\log n)$ some ladder pair of a long ladder has a short row cycle
through one marked column avoiding the other --- a lower-tail statement about rare local
configurations at rows selected by the square, whose heuristic expectation is
$\gtrsim\log^2n$.  The orbit method gives alternative sufficient conditions of the same
type, one of which trades the rarity for a genericity constant; and at the first ladder pair
alone it gives a third exact identity, expressing the bias $\Pr_X[B]-\Pr_X[A]$ as an
expectation over the event that none of the column pairs along the row permutation of that
pair admits a legal toggling turn, so that the conjecture also follows if that event has
probability $o(1/\log n)$ (in samples at $n\le100$ it is $\approx5.5/n$, which says nothing
about the required rate).  We state these precisely,
report the numerical evidence (up to $n=100$, including, through a sampler of completions of
a fixed rectangle, types of exponentially small probability), and locate the remaining
difficulty within the known switching technology: every exact identity we have found inside
the conditioned space is fair, and the lower bounds needed are of the kind whose proof in
\cite{CGW08} has a gap.  Along the
way we observe that the published proof of the upper bound in Lemma~3.12 of
Cavenagh--Greenhill--Wanless has a gap for every split other than $(2,2)$; our results do not
depend on that bound, and we re-derive from the surviving direction the tail estimate on
the number of cycles that the reduction theorem needs.
\end{abstract}

\section{Introduction}

Let $\PP_n$ be the law of the cycle type of $\sigma=\sigma_{1,2}(L)$ for a uniform Latin
square $L$ of order $n$, and $\QQ_n$ the law of the cycle type of a uniform derangement of
$[n]$.  Cavenagh, Greenhill and Wanless \cite[Conj.~6.1]{CGW08} conjectured:

\begin{conjecture}\label{conj:cgw}
$\dtv(\PP_n,\QQ_n)=o(1)$ as $n\to\infty$.
\end{conjecture}

For $\lambda\vdash n$ with all parts $\ge2$ let $\CC_n(\lambda)$ be the number of
completions of a fixed $2\times n$ Latin rectangle of cycle type $\lambda$ to a Latin
square, divided by $(n-2)!$ (the count does not depend on the rectangle chosen, since
rectangles of the same type are isotopic by a row-preserving isotopy), and let
$\gamma(\lambda)$ be the number of permutations of type $\lambda$.  Then
$\PP_n(\lambda)\propto\gamma(\lambda)\CC_n(\lambda)$ and
$\QQ_n(\lambda)\propto\gamma(\lambda)$, so the conjecture says that $\CC_n$ is nearly
constant on the types that carry most of the mass.  CGW proved
$\CC_n(\lambda)/\CC_n(\mu)\in[\tfrac12,\tfrac32]$ for adjacent types (one part of
$\lambda$ split into two parts of $\mu$, both $\ge2$), and used it to show that the two
laws agree within a factor $n^{3/2}$ on every type with at most $\tfrac65\log n$ parts; the upper bound $\tfrac32$, and
with it the upper half of the latter statement, is affected by the gap described in
Remark~\ref{rem:gap} and \cite{gapnote}.

\paragraph{Status of the arguments.}  The proofs below have been checked by automated
referees (independent language-model agents with access to the sources and the data) and,
where feasible, by exhaustive computation at $n\le7$ and by sampling (this includes
\S\ref{sec:orbit}(e), added last); no human referee has read them.

\paragraph{Results.}  Throughout, a \emph{mark} is a column pair $P=\{j,j'\}$ with
$j'=\sigma^\alpha(j)$ on an $m$-cycle of $\sigma$, $m=\alpha+\beta$, $\alpha,\beta\ge2$;
$X=X(n,\lambda;\alpha,\beta)$ is the set of marked squares of type $\lambda$; the
\emph{frame} of $P$ is the permutation of rows $\pi(r)=r'$ with $L(r',j')=L(r,j)$; and
$P$ is a \emph{$B$-pair} if rows $1,2$ lie in one cycle of $\pi$, an \emph{$A$-pair}
otherwise.  Our results are:
\begin{enumerate}[leftmargin=2em]
\item[(1)] \emph{Marked-pair identity} (Theorem~\ref{thm:marked}):
$2\CC_n(\lambda)/\CC_n(\mu)-1=\Pr_X[B]/\Pr_X[A]$ exactly, for every $n,\lambda,\alpha,\beta$.
\item[(2)] \emph{Reduction} (Theorem~\ref{thm:reduction}): if
$|\Pr_X[B]/\Pr_X[A]-1|\le\delta(n)$ for all adjacent pairs with $\alpha\ne\beta$
(equivalently $|\Pr_X[B]-\tfrac12|\le\tfrac14\delta(1+o(1))$), then
$\dtv(\PP_n,\QQ_n)=O(\delta\log n+n^{-1+o(1)})$.  The tail estimate it needs on the number
of cycles of $\sigma$ is proved from the trivial half of (1) alone
(Proposition~\ref{prop:tailsurvive}).
\item[(3)] \emph{Ladder identity} (Theorem~\ref{thm:ladder}): with the ladder pairs
$(x_k,y_k)=(\pi^{-k}(1),\pi^{-k}(2))$, $1\le k<K_0$, and ``flippable'' meaning that $j,j'$
lie in different cycles of the permutation $\rho_{x_k,y_k}$ of columns induced by rows $x_k,y_k$,
\[
\Pr_X[B]-\Pr_X[A]=\Pr_X[B,\mathrm{Flip}=\emptyset]-\Pr_X[A,\mathrm{Flip}=\emptyset],
\qquad |\Pr_X[B]-\tfrac12|\le\tfrac12\Pr_X[\mathrm{Flip}=\emptyset].
\]
So Conjecture~\ref{conj:cgw} follows from (L): $\Pr_X[\mathrm{Flip}=\emptyset]=o(1/\log n)$.
\item[(4)] \emph{Witness lemma} (Proposition~\ref{prop:Lwitness}): (L), hence the
conjecture, follows from Hypothesis~\ref{hyp:witness}: with probability $1-o(1/\log n)$,
a long ladder contains a pair $(x_k,y_k)$ whose $\rho_{x_k,y_k}$-cycle through $j$ has
length $\le\log^5n$ and avoids $j'$ (such a pair is flippable outright), and, for the offset
diagonals of short-arc instances, that a constant fraction of the expected number of
diagonals carry one.  This is the reference hypothesis of the paper.
\item[(5)] \emph{Orbit method} (Propositions~\ref{prop:orbit} and~\ref{prop:adaptive}):
the flippabilities of a family of ladder pairs are conditionally independent given the
orbit of a group of commuting column-cycle turns.  With a fixed matching of columns as
coordinates this gives $\Pr_X[\mathrm{Flip}_I=\emptyset]\le\E_X\exp(-\tfrac12\sum_{k\in I}\bar y_k)$
with an explicit ``orbit weight'' $\bar y_k$; with the \emph{adaptive} coordinates
$(\rho_k^t(j),\rho_k^t(j'))$ --- the $t$-th columns along the permutation $\rho_k=\rho_{x_k,y_k}$
of the $k$-th ladder pair from the marked columns --- every coordinate is separated by the
mark, so its turn toggles whenever it is legal and clean, and
$\Pr_X[\mathrm{Flip}_I=\emptyset]\le\E_X[2^{-|U|}]$, where $U$ is the set of ladder pairs
having a short clean candidate cycle (subject to a collision condition).  These give two
alternative sufficient conditions (Hypotheses~\ref{hyp:orbitgen} and~\ref{hyp:adaptive}),
for which we see no implication to or from Hypothesis~\ref{hyp:witness}: the first needs a
genericity constant that the switchings available inside $X$ cannot provide, the second
needs $C\log\log n$ rare local events on a long ladder where the witness lemma needs one.
\item[(6)] \emph{The first ladder pair alone} (Propositions~\ref{prop:firstpair}
and~\ref{prop:firstpairbias}): at the first pair $(x_1,y_1)$ the adaptive turn needs no
side conditions, and combined with the first ladder trade it gives
\[
\Pr_X[B]-\Pr_X[A]=\E_X\bigl[(\mathbf 1_B-\mathbf 1_A)(\mathbf 1_{F^c}-\mathbf 1_F);\,G^c\bigr],
\qquad\bigl|\Pr_X[B]-\tfrac12\bigr|\le\tfrac12\Pr_X[G^c],
\]
where $F$ is the flippability of $(x_1,y_1)$ and $G$ the event that one of the $\nu-1$
column pairs $\{\rho^t(j),\rho^t(j')\}$ along $\rho=\rho_{x_1,y_1}$ ($\nu$ = the length of
the shorter arc or cycle of $\rho$ at the mark) induces a permutation of rows that separates
$x_1$ from $y_1$ with a legal cycle.  So the conjecture also follows from
Hypothesis~\ref{hyp:firstpair}: $\sup_{\lambda,\alpha,\beta}\Pr_X[G^c]=o(1/\log n)$ --- an
event selected by one row pair, with candidates of constant probability (heuristically
$\Theta(n)$ of them; in samples of the type $(n)$ at $n=30,50,100$,
$\Pr_X[G^c]\approx5.5/n$) in place of a long ladder of rare events.  Again no implication
to or from Hypothesis~\ref{hyp:witness} is known.
\end{enumerate}
All four hypotheses are of one logical type: a lower bound on the probability of a
configuration at rows or columns selected by the frame of the mark, inside the space
conditioned on the type of rows $1,2$.  Their heuristic margins are large; none of the exact
tools of this paper gives a lower bound of that kind (\S\ref{sec:remains}).  The data
(Jacobson--Matthews samples, and a sampler of completions of a fixed $2\times n$ rectangle
that reaches types of $\PP_n$-probability $e^{-\Theta(n)}$, Remark~\ref{rem:fixrect}) give
$\Pr_X[B]=\tfrac12\pm0.003$ in every tested class at $n\le50$ and witness events that are
independent to three digits along the ladder; they do not reach the regime of any of the
hypotheses, and the empirical sizes quoted for the hypothesised quantities are first-moment
information only.

\paragraph{A gap in CGW's Lemma 3.12.}  Case 3 of the splitting procedure in
\cite{CGW08} (cross-switch at $\{\omega j,\omega j'\}$, then backflip at $\{j,j'\}$) does
not behave as claimed: the cross-switch reverses the $\omega$-arc from $\omega j$ to
$\omega j'$ (or the complementary arc, depending on the smallest corner symbol), which
contains $j'$ (resp.\ $j$), so $\{j,j'\}$ ends at distance $2$.  For $\min(\alpha,\beta)\ge3$
the backflip then lands in the wrong type; for $\min(\alpha,\beta)=2<\max(\alpha,\beta)$ the
type is right but the edge is not an edge of the joining multigraph, so the identity
$G_S=G_J$ on which the double count rests fails (already at $n=5$, where we computed both
multigraphs exactly).  We give the details in a separate note \cite{gapnote}; here we only record
(Remark~\ref{rem:gap}) which published statements depend on it.  Nothing in the present
paper does: the lower bound $\tfrac12$ of CGW's lemma is the trivial half of
Theorem~\ref{thm:marked}, and the upper bound $\tfrac32$ is, by the same theorem,
$\Pr_X[B]\le\tfrac23$ --- a weak form of what we are after.  By Theorem~\ref{thm:ladder}
it would follow from $\Pr_X[\mathrm{Flip}=\emptyset]\le\tfrac13$.

\paragraph{Conventions.}  $L(r,c)$ is the entry in row $r$, column $c$.  For rows $x,y$,
$\rho_{x,y}$ is the permutation of columns with $L(y,\rho_{x,y}(q))=L(x,q)$, so that
$\sigma=\rho_{1,2}$ up to inversion; a \emph{row cycle} of $(x,y)$ is a cycle of
$\rho_{x,y}$.  For columns $q,c$, $\delta_{q,c}$ is the permutation of rows with
$L(\delta(r),c)=L(r,q)$, and its cycles are the \emph{column cycles} of $(q,c)$.  A
\emph{trade} of rows $x,y$ along a row cycle swaps their entries on the columns of that
cycle; a \emph{turn} of a column cycle swaps the entries of columns $q,c$ in its rows.  Both
are Latin trades (the result is again a Latin square).  For $\lambda\vdash n$, $S(\lambda)$ is
the set of Latin squares of order $n$ with first row the identity whose rows $1,2$ have
cycle type $\lambda$, so $|S(\lambda)|=\gamma(\lambda)(n-2)!\,\CC_n(\lambda)$; we also
write $D_\lambda=\gamma(\lambda)$ for the number of derangements of type $\lambda$.  A
column pair $\{j,j'\}$ is \emph{admissible} if $j'\notin\{j,\sigma(j),\sigma^{-1}(j)\}$
(CGW's condition for the pair operations; every marked pair is admissible since
$\alpha,\beta\ge2$).  Following CGW, the \emph{flip} at an $A$-pair whose columns lie in
different $\sigma$-cycles, and the \emph{backflip} (or \emph{unflip}) at an $A$-pair whose
columns lie in one $\sigma$-cycle, is the turn of the column cycle of the pair through row
$2$ (it joins, resp.\ splits, the $\sigma$-cycles involved); the \emph{switch} at a pair
whose columns lie in different $\sigma$-cycles is the trade of rows $1,2$ along the row
cycle through one of the two columns, selected by CGW's rule (the cycle whose columns
carry the smaller minimum symbol).  We call Conjecture~\ref{conj:cgw}, which asks only
for convergence in total variation, the \emph{weak} conjecture; the stronger statements
suggested by Cameron (``tends very rapidly'') and by the data (pointwise ratios
$1+o(1)$) are not addressed here.
Jacobson--Matthews sampling
\cite{JM96} is used for all numerical data; the scripts and logs are in the project
repository.

\section{The marked-pair identity}\label{sec:marked}

Throughout, squares are normalised to have first row $\mathrm{id}$; CGW's column map
$\omega$, defined by $L(1,\omega(j))=L(2,j)$, is then $\omega(j)=L(2,j)$, and we take
$\sigma=\omega$ (only the cycle type matters, and it is the same for $\omega$ and $\omega^{-1}$).  For the proofs it
is convenient to work with \emph{unnormalised} squares (equivalently,
symbol-relabelling classes): every instance count below then acquires a
common factor $n!$, which cancels from every ratio, and the maps we use
(flip, switch) act on unnormalised squares without any renormalisation
step.  Statuses, cycle types, and column positions are
relabelling-invariant, so we freely state results in the normalised
picture.

\begin{definition}[marked instances]\label{def:marked}
Let $\lambda\vdash n$ contain a part $m=\alpha+\beta$
($\alpha,\beta\ge2$), and let $\mu$ be $\lambda$ with one part $m$
split into $\alpha,\beta$.  The \emph{splitting instance space} is
\[
X \;=\; X(n,\lambda;\alpha,\beta)
\;=\;\bigl\{(L',P)\;:\;L'\in S(\lambda),\;
P=\{j,\sigma^\alpha(j)\}\text{ with } j\text{ in an $m$-cycle of }
\sigma(L')\bigr\},
\]
unordered pairs counted once.  Thus every $L'$ carries exactly
$\lambda_m\, m$ marked pairs when $\alpha\neq\beta$ and
$\lambda_m\,\alpha$ when $\alpha=\beta$.  Since $\alpha,\beta\ge 2$
every marked pair is admissible, and $X=X_A\sqcup X_B$ according to the
$A$/$B$ status of $P$ in $L'$.  The \emph{joining instance space} is
\[
J \;=\; J(n,\mu;\alpha,\beta)
\;=\;\bigl\{(L,P)\;:\;L\in S(\mu),\;P=\{j,j'\},\;
j\text{ in an $\alpha$-cycle},\; j'\text{ in a different
$\beta$-cycle}\bigr\},
\]
again with $J=J_A\sqcup J_B$ by the status of $P$ in $L$.
\end{definition}

The counting arithmetic is immediate: with
$D_\nu$ the number of derangements of cycle type $\nu$ and
$N(\nu)=\CC_n(\nu)(n-2)!$,
\begin{equation}\label{eq:Xarith}
\frac{2|X|}{|J|}
=\frac{2\,D_\lambda N(\lambda)\,\lambda_m m}
       {D_\mu N(\mu)\,\mu_\alpha\mu_\beta\,\alpha\beta}
=\frac{2\CC_n(\lambda)}{\CC_n(\mu)}
\qquad(\alpha\neq\beta),
\end{equation}
using $D_\lambda/D_\mu=\alpha\mu_\alpha\,\beta\mu_\beta/(m\lambda_m)$;
for $\alpha=\beta$ one has
$D_\lambda/D_\mu=\alpha^2\mu_\alpha(\mu_\alpha-1)/(m\lambda_m)$,
$|J|=|S(\mu)|\binom{\mu_\alpha}{2}\alpha^2$ and
$|X|=|S(\lambda)|\lambda_m\alpha$, and \eqref{eq:Xarith} holds
verbatim.  The content of the identity is therefore the following
count.

\begin{lemma}[cycle trades preserve pair status]\label{lem:cycletrade}
Let $P=\{j,j'\}$ be an admissible column pair of a Latin square $L$ and
let $C$ be one column cycle of $P$.  Trading $C$ (swapping the entries
of columns $j,j'$ in the rows of $C$) inverts the pair permutation
$\pi_P$ on $C$ and fixes it elsewhere; the partition of rows into
column cycles of $P$ is unchanged.  In particular the $A$/$B$ status of
$P$ itself is invariant under the CGW flip at $P$, which trades the
full column cycle through row~$2$.
\end{lemma}

\begin{proof}
Write $\pi=\pi_P$, so $\pi(r)$ is the row of column $j'$ holding the
symbol $L(r,j)$; $C$ is $\pi$-closed and the trade swaps
$L(r,j)\leftrightarrow L(r,j')$ for $r\in C$.  Fix $r\in C$ and let
$\pi'$ be the new pair permutation.  The new $L(r,j)$ is the old
$L(r,j')$; the row of the new column $j'$ holding that symbol is found
by solving (new $L(s,j')$) $=$ (old $L(r,j')$): for $s\in C$ the new
$L(s,j')$ is the old $L(s,j)$, and old $L(s,j)=$ old $L(r,j')$ exactly
when $\pi(s)=r$.  Hence $\pi'=\pi^{-1}$ on $C$; off $C$ nothing
changes.  A cycle inverted in place has the same support, so the
partition into cycles is unchanged.
\end{proof}

\begin{lemma}[switch toggle]\label{lem:switchtoggle}
Let $(L,P)\in J$ and let $Q$ be the row cycle through one column of
$P$ (the two row cycles are distinct), selected by CGW's min-symbol rule.  Trading $Q$ (the CGW switch)
changes $\pi_P$ by precomposition with the transposition of rows
$1,2$: the cycles of $\pi_P$ through rows $1,2$ are split if rows
$1,2$ are in one cycle, and joined if in two.  Hence the switch
toggles the status of $P$, preserves $S(\mu)$ and the instance data
(it inverts one $\sigma$-cycle in place), and is an involution:
the trade preserves the traded cycle's symbol set (the two rows carry
the same symbols on its columns), so the min-symbol rule re-selects the
same cycle.  Consequently $|J_A|=|J_B|$, exactly.
\end{lemma}

\begin{proof}
Say $Q$ is the row cycle through $j$; its columns include $j$ but not
$j'$.  The trade swaps the row-1 and row-2 entries in the columns of
$Q$; for $\pi_P$ only column $j$ matters, where the entries of rows
$1,2$ are exchanged.  Thus $\pi_P'=\pi_P\circ\tau_{12}$, and
multiplying by a transposition splits or joins the cycles through its
two points.  Rows $1,2$ lie in one cycle of $\pi_P$ exactly when $P$
is a $B$-pair.  The trade swaps the two rows' entries along the
columns of $Q$, which inverts the corresponding $\sigma$-cycle in
place: the cycle type and the column sets of all row cycles are
preserved, so $(L,P)\mapsto(\mathrm{sw}(L),P)$ maps $J$ to $J$.
\end{proof}

\begin{theorem}[marked-pair identity]\label{thm:marked}
For every $n$, $\lambda$, and split $m=\alpha+\beta$ with
$\alpha,\beta\ge2$ (including $\alpha=\beta$),
\[
|J_A|\;=\;|J_B|\;=\;|X_A|,
\qquad\text{hence}\qquad
\boxed{\;\frac{\CC_n(\lambda)}{\CC_n(\mu)}\;=\;\frac{|X|}{2|X_A|},
\qquad
2\,\frac{\CC_n(\lambda)}{\CC_n(\mu)}-1\;=\;\frac{|X_B|}{|X_A|}.\;}
\]
In particular $\CC_n(\lambda)>\tfrac12\,\CC_n(\mu)$ unconditionally,
and Hypothesis~\ref{hyp:eq} ($\mathrm{EQ}(\delta)$) is equivalent to
\[
\Bigl|\;\Pr\bigl(P\text{ is a $B$-pair}\bigr)-\tfrac12\;\Bigr|
\;\le\;\tfrac{\delta}{4}(1+o(1)),
\]
for $(L',P)$ uniform on $X$: the law of a single bit --- whether rows
$1,2$ lie in one column cycle of the marked pair
$\{j,\sigma^\alpha(j)\}$ --- of a single uniform square of
$S(\lambda)$.
\end{theorem}

\begin{proof}
By \eqref{eq:Xarith} it suffices to prove $|J|=2|X_A|$.
The flip at an $A$-pair $P$ whose columns lie in one row cycle at
$\sigma$-distances $(\alpha,\beta)$ splits that cycle into an
$\alpha$- and a $\beta$-cycle and maps $S(\lambda)\to S(\mu)$; by
Lemma~\ref{lem:cycletrade} the pair $P$ is still an $A$-pair
afterwards, so $x\mapsto\mathrm{unflip}(x)$ maps $X_A$ into $J_A$.
The flip at a fixed pair is an involution (the traded cycle again
contains row 2, with the same row support), so
$\mathrm{join}=\mathrm{flip}$ inverts it: $X_A\leftrightarrow J_A$ is
a bijection.  Composing with the switch of
Lemma~\ref{lem:switchtoggle} gives a bijection
$X_A\leftrightarrow J_B$, namely
$x\mapsto \mathrm{sw}(\mathrm{unflip}(x))$ with inverse
$(L,P)\mapsto \mathrm{flip}(\mathrm{sw}(L))$ --- the map
``switch-then-flip'', which is exactly the CGW joining step at a
$B$-pair.  Hence $|J|=|J_A|+|J_B|=2|X_A|$.
\end{proof}


\section{The reduction theorem}\label{sec:reduction}

\begin{hypothesis}[B-pair equidistribution, $\mathrm{EQ}(\delta)$]\label{hyp:eq}
For every $\lambda\vdash n$ with a part $m=\alpha+\beta$, $\alpha\ne\beta$, $\alpha,\beta\ge2$, and $\mu$ the split type,
$\bigl|\bar q(\mu\to\lambda)-1\bigr|\le\delta(n)$, where $\bar q(\mu\to\lambda):=2\CC_n(\lambda)/\CC_n(\mu)-1=\Pr_X[B]/\Pr_X[A]$ by Theorem~\ref{thm:marked}.
\end{hypothesis}

By Theorem~\ref{thm:marked}, $\mathrm{EQ}(\delta)$ is the statement $|\Pr_X[B]-\tfrac12|\le\tfrac\delta4(1+o(1))$ for all adjacent pairs with $\alpha\ne\beta$.  The proof below needs a tail bound on the number $\kappa$ of cycles of $\sigma$; we first prove it from the trivial half of Theorem~\ref{thm:marked} alone, so that nothing depends on the upper bound of \cite[Lemma~3.12]{CGW08} (Remark~\ref{rem:gap}).

\begin{proposition}[the number of cycles, from the surviving direction only]\label{prop:tailsurvive}
Let $L$ be a uniform Latin square of order $n$, $\sigma$ the column
permutation of rows $1,2$, $N_\ell$ its number of $\ell$-cycles and
$\kappa=\sum_\ell N_\ell$.  Then, using only the identity
$|J|=2|X_A|$ of Theorem~\ref{thm:marked} and $|X_A|\le|X|$:
\begin{enumerate}[leftmargin=2em]
\item[(i)] $\E[\ell N_\ell]\le2n/(n-\ell)$ for every $\ell<n$; in
particular $\E[N_\ell]\le4/\ell$ for $\ell\le n/2$.
\item[(ii)] For every $k\ge1$, with $(\kappa)_k=\kappa(\kappa-1)\cdots(\kappa-k+1)$,
\[
\E\bigl[(\kappa)_k\bigr]\;\le\;2^{k-1}\sum_{\ell_1,\dots,\ell_k\ge2}\frac{\E[mN_m]}{\ell_1\cdots\ell_k}\Big|_{m=\ell_1+\cdots+\ell_k}
\;\le\;2^{k-1}\Bigl(4+\frac{k^2}{\log n}\Bigr)(\log n)^k .
\]
\item[(iii)] For $A\ge3$, $\Pr[\kappa\ge A\log n]\le(7+\log n)\,n^{-\log((A-1)/2)}$,
which is $O(n^{-2}\log n)$ for $A\ge16$;
and $\Pr[\text{all cycles of }\sigma\text{ have the same length }\ell<n]\le4d(n)/n$,
$d(n)$ the number of divisors.
\end{enumerate}
\end{proposition}

\begin{proof}
(i) For $\beta\ne\ell$ the joining instances $(L,\{j,j'\})$ with $j$ in
an $\ell$-cycle and $j'$ in a $\beta$-cycle number $\ell N_\ell\beta N_\beta$
per square, the splitting instances (pairs at distance $\ell$ on an
$(\ell+\beta)$-cycle) number $(\ell+\beta)N_{\ell+\beta}$, and
$|J|=2|X_A|\le2|X|$ gives
$\E[\ell N_\ell\,\beta N_\beta]\le2\E[(\ell+\beta)N_{\ell+\beta}]$; for
$\beta=\ell$ the same count with $\ell^2\binom{N_\ell}{2}$ and
$\ell N_{2\ell}$ gives $\E[\ell^2N_\ell(N_\ell-1)]\le4\ell\E[N_{2\ell}]$.
Summing the first over $\beta\ne\ell$ ($\sum_\beta\beta N_\beta=n$) and
adding the second (its right side is the missing summand $m=2\ell$):
$\E[\ell N_\ell(n-\ell N_\ell)]+\E[\ell^2N_\ell(N_\ell-1)]\le2\E[\sum_{m\ge\ell+2}mN_m]\le2n$.  Hence
$n\E[\ell N_\ell]=\E[\ell N_\ell(n-\ell N_\ell)]+\E[\ell^2N_\ell(N_\ell-1)]+\ell\E[\ell N_\ell]
\le2n+\ell\E[\ell N_\ell]$.

(ii) Count the instances $(L,(C_1,j_1),\dots,(C_k,j_k))$: $C_i$
pairwise distinct $\sigma$-cycles of $L$ of lengths $\ell_i$ and
$j_i\in C_i$; there are $\ell_1\cdots\ell_k$ of them per ordered
$k$-tuple of distinct cycles with these lengths.  Join $C_1$ and $C_2$
at the pair $\{j_1,j_2\}$ by a joining step: flip if the pair is an
$A$-pair; if it is a $B$-pair, first trade the row cycle $C_2$ of $j_2$
(a switch in the sense of Lemma~\ref{lem:switchtoggle}, whose proof
applies to the row cycle through either column of the pair; it inverts
$C_2$ in place and does not touch the cycle of $j_1$), then flip.  The
flip replaces $\sigma$ by $\sigma\circ(j_1\,j_2)$ on rows $1,2$, so
the result has $j_1,j_2$ on one cycle of length $\ell_1+\ell_2$ with
$j_2=\sigma'^{\ell_2}(j_1)$ whether or not the trade was applied
(CGW Lemma~3.3), every other cycle of $\sigma$ (in particular
$C_3,\dots,C_k$) untouched, and the map is at most two-to-one (the
backflip at $\{j_1,j_2\}$ returns $L$ or the traded square).  Repeat
with $\{j_1,j_3\}$, \dots, $\{j_1,j_k\}$, always trading the cycle of
the new column when needed, so that the cycle containing $j_1$ is never
inverted: after step $i$ every earlier $j_r$ is shifted by exactly
$\ell_{i+1}$.  After $k-1$ steps all $j_i$ lie on one cycle of length
$m=\sum\ell_i$ at the positions $j_i=\sigma'^{\ell_i+\cdots+\ell_k}(j_1)$,
determined by $j_1$ and $(\ell_1,\dots,\ell_k)$, and the composite is at
most $2^{k-1}$-to-one.  (With CGW's min-symbol rule the trade may
invert the merged cycle and the positions are not determined; the
count would then be $2^{2k-3}$-to-one.)  Hence the number of instances with lengths
$(\ell_i)$ is at most $2^{k-1}\sum_{L'}mN_m(L')$, i.e.\
$\E[\ell_1\cdots\ell_k\,T(\ell_1,\dots,\ell_k)]\le2^{k-1}\E[mN_m]$ with
$T$ the number of ordered $k$-tuples of distinct cycles with these
lengths; summing over $(\ell_i)$ gives the first bound, since
$\sum_{(\ell_i)}T=(\kappa)_k$.  For the second, split at $m\le n/2$,
where $\E[mN_m]\le4$ by (i) and
$\sum_{m\le n/2}\sum_{\sum\ell_i=m}\prod\ell_i^{-1}\le(\sum_{\ell\le n/2}\ell^{-1})^k\le(\log n)^k$,
and $m>n/2$, where some $\ell_i\ge m/k$, so
$\sum_{\sum\ell_i=m}\prod\ell_i^{-1}\le k\cdot(k/m)(\log n)^{k-1}$, and
$\sum_{m>n/2}mN_m\cdot k^2(\log n)^{k-1}/m\le k^2(\log n)^{k-1}$ because at
most one cycle is longer than $n/2$.

(iii) Markov for $(\kappa)_k$ with $k=\lceil\log n\rceil\le\log n+1$:
on $\{\kappa\ge K\}$, $(\kappa)_k\ge(K-k+1)^k\ge((A-1)\log n)^k$ for
$K=A\log n$, and $k^2/\log n\le\log n+3$, so
$\Pr[\kappa\ge K]\le2^{k-1}(4+k^2/\log n)(\log n)^k/((A-1)\log n)^k
\le(7+\log n)(2/(A-1))^{k}\le(7+\log n)(2/(A-1))^{\log n}$ for $A\ge3$,
and $(A-1)/2\ge e^2$ for $A\ge16$.  For the equal-length types,
$\ell\mid n$ and $\ell<n$ force $\ell\le n/2$, so
$\Pr[N_\ell=n/\ell]\le\E[N_\ell]\ell/n\le4/n$; $\ell=n$ is the $n$-cycle.
\end{proof}

\begin{theorem}\label{thm:reduction}
Assume $\mathrm{EQ}(\delta)$ with $\delta=\delta(n)\to0$.
Then
\[
\dtv(\PP_n,\QQ_n)\;=\;O\bigl(\delta\,\log n\;+\;n^{-1+o(1)}\bigr).
\]
In particular $\mathrm{EQ}(o(1/\log n))$ implies
Conjecture~\ref{conj:cgw}.
\end{theorem}

\begin{proof}
Call $\lambda$ \emph{plain} if not all its parts are equal (so the only
non-plain types are $(\ell^{\kappa})$, $\kappa\ge2$).

\emph{Routing.}  Every plain $\lambda$ with $\kappa=\kappa(\lambda)$ parts is
reachable from $(n)$ by a chain of $\kappa-1$ splits, each with
$\alpha\ne\beta$.  Build the chain in reverse, splitting parts of $\lambda$
off a shrinking residual, in a suitable order; a collision means splitting a
part $p$ off a residual $R$ with $p=R-p$.  If the two largest parts of
$\lambda$ are unequal, split smallest-first and let them be the final pair:
while at least three parts remain, the split-off minimum $p$ satisfies
$R-p\ge2p>p$ (all parts are $\ge2$ and at least two other parts remain), so
no collision occurs, and the final split is unequal by hypothesis.  If
instead the two largest parts are equal, say to $\ell_{\max}$, then
$\ell_{\max}\neq n/2$ (else $\lambda=(n/2,n/2)$ is not plain), so we may
split one copy of $\ell_{\max}$ off first without collision, then proceed
smallest-first among the remaining parts, reserving for the final split the
pair $(b,\ell_{\max})$ where $b<\ell_{\max}$ is the smallest part: at every
intermediate step the residual minus the split-off minimum still contains
the reserved $\ell_{\max}$ and the reserved $b$, hence exceeds the
minimum (also when a further copy of $\ell_{\max}$ is the split-off
minimum: $R-p\ge\ell_{\max}+b>p$), and the final split is unequal.

For each step of such a chain, Theorem~\ref{thm:marked} and
$\mathrm{EQ}(\delta)$ give
$|S(p_{i+1})|/|S(p_i)|=(\gamma(p_{i+1})/\gamma(p_i))(1+\varepsilon_i)$ with
$|\varepsilon_i|\le\delta$.  Multiplying along the chain,
\[
\frac{\PP_n(\lambda)}{\PP_n((n))}
=\frac{\gamma(\lambda)}{\gamma((n))}\prod_{i=1}^{\kappa-1}(1+\varepsilon_i)^{\pm1},
\qquad |\varepsilon_i|\le\delta,
\]
so for all plain $\lambda$ with $\kappa(\lambda)\le K$,
$\PP_n(\lambda)=\QQ_n(\lambda)\,Z^{-1}(1+O(K\delta))$ (the bound is vacuous unless $K\delta=O(1)$, which is the case of interest) where $Z>0$ is a
normalising constant independent of $\lambda$ (namely
$Z=\sum_\nu\gamma(\nu)\CC_n(\nu)/(|\DD_n|\,\CC_n((n)))$).

Let $K=\lceil A\log n\rceil$ with $A$ a large absolute constant, and let
$\mathcal B$ be the set of types that are either non-plain or have
$\kappa>K$.  Under $\QQ_n$, Markov's inequality with the generating function
$\sum_{\pi\in\DD_n}2^{\kappa(\pi)}=n!\,[x^n]e^{-2x}(1-x)^{-2}\le n\cdot n!$
(cf.\ \cite[proof of Lemma~4.4]{CGW08}) gives
$\QQ_n(\kappa>K)\le e\,n\,2^{-K}=O(n^{-2})$ for $A\log2\ge3$; also
$\QQ_n(\text{non-plain})\le\sum_{\ell\mid n,\ \ell<n}
\gamma((\ell^{n/\ell}))/D_n=O(n^{-1})$.  Under $\PP_n$, Proposition~\ref{prop:tailsurvive} gives
$\PP_n(\kappa>K)=O(n^{-2}\log n)$ for $A\ge16$ and
$\PP_n(\text{non-plain})=O(n^{-1+o(1)})$ without any input from
\cite{CGW08} beyond the identity of Theorem~\ref{thm:marked}.

On the complement of $\mathcal B$, summing
$\PP_n=\QQ_n Z^{-1}(1+O(K\delta))$ over $\lambda\notin\mathcal B$ and using
the tail bounds shows $Z^{-1}=1+O(K\delta)+O(n^{-1+o(1)})$, whence
$\sum_{\lambda\notin\mathcal B}|\PP_n-\QQ_n|=O(K\delta)+O(n^{-1+o(1)})$.
Combining with the tails,
$\dtv(\PP_n,\QQ_n)=O(\delta\log n+n^{-1+o(1)})$.
\end{proof}

\begin{remark}[the gap in CGW's Lemma 3.12]\label{rem:gap}
In \cite{CGW08}, the direction $|S(\lambda,F)|N_\lambda\ge\tfrac12|S(\mu,F)|N_\mu$ of
Lemma~3.12 uses only that each splitting pair carries at most two edges and each joining
pair at least one; by Theorem~\ref{thm:marked} it is the trivial inequality $|X|\ge|X_A|$.
The direction $\le\tfrac32$ requires every splitting pair of every $\lambda$-square to carry
exactly two edges that are also edges of the joining multigraph, and case~3 of the splitting
procedure fails to supply them for every split other than $(2,2)$ \cite{gapnote}.
Consequently the upper bound of CGW's Theorem~3.13, the lower bound of their Lemma~4.1,
the parity bounds of Theorem~4.3, the bound $\Pr[\sigma\text{ is an $n$-cycle}]\le2n^{-2/3}$
of Lemma~4.4, both halves of Corollary~4.5, and the results of their Section~4.2 that rest
on these (Corollary~4.6, Lemma~4.7, Theorems~4.9 and~4.11, Corollary~4.10 --- the last
superseded by \cite{KS18}) are not established by the published argument.  For
$\min(\alpha,\beta)=2$ the count can be corrected with the constant $2$ in place of $\tfrac32$
\cite[Prop.~3]{gapnote}: for marks at distance $2$ or $m-2$ on an $m$-cycle,
$\CC_n(\lambda)/\CC_n(\mu)\le2$, i.e.\ $\Pr_X[A]\ge\tfrac14$ unconditionally.  The present
paper uses none of them:
Proposition~\ref{prop:tailsurvive} replaces Corollary~4.5 in the proof of
Theorem~\ref{thm:reduction}.  The exact completion counts tabulated in \cite{CGW08} give
$\CC_n(\lambda)/\CC_n(\mu)\in[0.985,1.0035]$ for every adjacent pair at $7\le n\le11$ (and
$=\tfrac32$ at $n=5$, $\tfrac78,\tfrac{11}{14},\tfrac{21}{22}$ at $n=6$), so the statement
itself is not in doubt.
\end{remark}

\section{The ladder identity}\label{sec:ladder}

Notation:
$(L,P)\in X=X(n,\lambda;\alpha,\beta)$, $P=\{j,j'\}$, frame
$\pi=\pi_P$ ($L(\pi(r),j')=L(r,j)$), $C_r$ the $\pi$-cycle of row $r$,
$\rho_{x,y}$ the column permutation of rows $x,y$
($L(y,\rho_{x,y}(q))=L(x,q)$), $(x,y)$ \emph{flippable} iff $j,j'$ lie
in different cycles of $\rho_{x,y}$, and the \emph{$j'$-trade} at
$(x,y)$ exchanges rows $x,y$ on the $\rho_{x,y}$-cycle through $j'$,
which replaces $\pi$ by $(x\,y)\circ\pi$ if $(x,y)$ is flippable and by
$(x\,y)\circ\pi\circ(x\,y)$ otherwise.  $B=\{1\sim2\}$ (rows $1,2$ in a
common frame cycle), $A$ its complement, $\mathbf 1_B$ the bit; for $B$-instances $d_{12}$ and
$d_{21}$ are the numbers of $\pi$-steps from $1$ to $2$ and from $2$ to
$1$.

\begin{definition}[ladder]\label{def:ladder}
Put $K_0=\min(|C_1|,|C_2|)$ on $A$ and $K_0=\min(d_{12},d_{21})$ on
$B$, and for $1\le k<K_0$ let
\[
x_k=\pi^{-k}(1),\qquad y_k=\pi^{-k}(2),\qquad
T_k=\text{the $j'$-trade at }(x_k,y_k).
\]
The pairs $\{x_k,y_k\}$ are the \emph{ladder} of $(L,P)$, and
$\mathrm{Flip}(L,P)=\{k<K_0:(x_k,y_k)\text{ flippable}\}$.
\end{definition}

The rows $x_1,\dots,x_{K_0-1}$ are distinct rows of $C_1\setminus\{1\}$
(on $A$) or of the arc from $2$ to $1$ (on $B$), the $y_k$ likewise in
$C_2\setminus\{2\}$ or the arc from $1$ to $2$; so the $2(K_0-1)$
ladder rows are distinct and avoid rows $1,2$, and $K_0\ge2$: on $A$
because $\pi$ is a derangement, on $B$ by the admissibility of the mark
($\pi(1)\ne2$, $\pi(2)\ne1$).

\begin{theorem}[ladder identity]\label{thm:ladder}
For every $n$, $\lambda$ and split $\alpha+\beta=m$ with
$\alpha,\beta\ge2$,
\[
\Pr_X[B]-\Pr_X[A]\;=\;\Pr_X\bigl[B,\ \mathrm{Flip}=\emptyset\bigr]
-\Pr_X\bigl[A,\ \mathrm{Flip}=\emptyset\bigr],
\qquad\text{hence}\qquad
\bigl|\Pr_X[B]-\tfrac12\bigr|\;\le\;\tfrac12\Pr_X[\mathrm{Flip}=\emptyset].
\]
The same identity holds with $X$ replaced by the set of squares with a fixed pair of
rows $1,2$ (of type $\lambda$) and a fixed mark $P$.
\end{theorem}

\begin{proof}
Fix $(L,P)\in X$ and write $K=\{1,\dots,K_0-1\}$.  The ladder depends
on $(L,P)$, so the $T_k$ are not maps of $X$; the statement to prove is
that $(L,P)\sim T_S(L,P)$, $S\subseteq K$ (the trades at the
unordered pairs of $S$, in any order), is an equivalence relation on
$X$ whose classes have exactly $2^{|K|}$ elements and on which the bit
is $\mathbf 1_B(T_S(L,P))=\mathbf 1_B(L,P)+|S\cap\mathrm{Flip}|\bmod2$.
This follows once we know that each $T_k$ is an involution of $X$
which preserves $K_0$, the set of unordered ladder pairs and the
flippability of each of them, that $T_k$ toggles the bit exactly when
$k\in\mathrm{Flip}$, and that trades at disjoint row pairs commute:
then the same trades, at the same unordered pairs with the same
flippabilities, are available from every point of a class (symmetry
and transitivity), the $2^{|K|}$ points $T_S(L,P)$ are distinct
(rows $x_k,y_k$ are changed by $T_k$ and by no other $T_i$), and the
bit formula holds.  If $\mathrm{Flip}\ne\emptyset$ exactly half of the
subsets $S$ have $|S\cap\mathrm{Flip}|$ odd, so $B$ and $A$ each
contain half of the class; summing over classes gives the identity,
and the bound follows from $\Pr[B]+\Pr[A]=1$.

\emph{$T_k$ is a well-defined involution of $X$.}  The $j'$-trade at
$(x_k,y_k)$ changes only rows $x_k,y_k\notin\{1,2\}$, so $\sigma$ and
the mark are untouched and $(L,P)$ stays in $X$; it is an involution
because the traded columns are again the $\rho'_{x_k,y_k}$-cycle
through $j'$.

\emph{Effect on the frame.}  If $k\in\mathrm{Flip}$ the new frame is
$\pi'=(x_k\,y_k)\circ\pi$.  On $A$ this merges $C_1$ and $C_2$: the
new cycle runs $1\to\pi(1)\to\dots\to x_{k+1}\to y_k\to\pi(y_k)\to\dots
\to2\to\dots\to y_{k+1}\to x_k\to\dots\to1$, so $d'_{12}=|C_1|$ and
$d'_{21}=|C_2|$ (reading $x_{|C_1|}=1$, $y_{|C_2|}=2$).  On $B$,
$x_k$ lies on the arc from $2$ to $1$ and $y_k$ on the arc from $1$ to
$2$ (as $k<\min(d_{12},d_{21})$), and the split into
$(x_k,\pi(x_k),\dots,\pi^{-1}(y_k))$ and $(y_k,\dots,\pi^{-1}(x_k))$
puts $1$ ($k$ steps after $x_k$, fewer than the $d_{12}$ steps from
$x_k$ to $y_k$) in the first cycle and $2$ in the second: the new
state is in $A$ with $|C'_1|=d_{12}$, $|C'_2|=d_{21}$.  In both cases
$K_0$ is preserved and the bit toggles.  If $k\notin\mathrm{Flip}$ the
new frame is $(x_k\,y_k)\circ\pi\circ(x_k\,y_k)$: the cycle partition
is unchanged up to exchanging the labels $x_k,y_k$, so $K_0$ and the
bit are preserved.

\emph{The ladder is preserved.}  Let $\pi'=(x_k\,y_k)\circ\pi$ (the
conjugation case is similar and easier).  Since
$\pi'^{-1}(r)=\pi^{-1}((x_k\,y_k)(r))$, the preimage chains of $1$ and
$2$ under $\pi'$ agree with those under $\pi$ up to depth $k$ and are
exchanged beyond it: $\pi'^{-i}(1)=x_i$ and $\pi'^{-i}(2)=y_i$ for
$i\le k$, while $\pi'^{-(k+1)}(1)=\pi^{-1}(y_k)=y_{k+1}$ and then
$\pi'^{-i}(1)=y_i$, $\pi'^{-i}(2)=x_i$ for $k<i<K_0$.  So the unordered
pairs $\{x_i,y_i\}$, $i\in K$, are the ladder of $T_k(L,P)$.  In the
conjugation case the same computation with $(x_k\,y_k)\circ\pi\circ(x_k\,y_k)$
gives $\pi'^{-i}(1)=(x_k\,y_k)(x_i)$, which again permutes the pairs
$\{x_i,y_i\}$ among themselves (only $i=k$ is affected, by exchanging
$x_k,y_k$).

\emph{Flippability is preserved and the $T_k$ commute.}  $T_k$ changes
rows $x_k,y_k$ only, and the flippability of $(x_i,y_i)$, $i\ne k$,
depends only on rows $x_i,y_i$, which are different rows.  Likewise
$T_i$ and $T_k$ act on disjoint pairs of rows, each along a column set
determined by its own two rows, so $T_iT_k=T_kT_i$.  (After $T_k$ the
pair $\{x_i,y_i\}$ may have changed its index, but the trade at an
unordered pair does not depend on the index or on the order of the
two rows.)

The statement for fixed $R$ holds because every $T_k$ fixes rows $1$
and $2$.
\end{proof}

\begin{corollary}[the conjecture reduces to the ladder]\label{cor:ladder}
Write $f(n,\lambda;\alpha,\beta)=\Pr_X[\mathrm{Flip}=\emptyset]$.  By
Theorem~\ref{thm:marked}, $\bar q(\mu\to\lambda)-1=(\Pr_X[B]-\Pr_X[A])/\Pr_X[A]$
and $\Pr_X[A]\ge(1-f)/2$, so $|\bar q-1|\le2f/(1-f)$: $\mathrm{EQ}(\delta)$
holds with $\delta(n)=\max_{\lambda,\alpha\ne\beta}2f/(1-f)$.  If this
$\delta(n)\to0$, Theorem~\ref{thm:reduction} gives
\[
\dtv(\PP_n,\QQ_n)=O\Bigl(\log n\cdot\max_{\lambda,\alpha\ne\beta}
f(n,\lambda;\alpha,\beta)+n^{-1+o(1)}\Bigr).
\]
In particular the weak CGW conjecture follows from
\begin{equation}\label{eq:L}
\textup{(L)}\qquad \Pr_X\bigl[\text{no ladder pair is flippable}\bigr]=o(1/\log n)
\quad\text{uniformly in }\lambda,\alpha,\beta .
\end{equation}
\end{corollary}


\begin{remark}[data]\label{rem:ladderdata}
On Jacobson--Matthews samples with a uniformly random mark ($9000$, $7200$, $3000$ ladder-pair
instances at $n=30,50,100$; \texttt{runs/s13/f1/}): $\Pr[\mathrm{Flip}=\emptyset\mid K_0]=2^{-(K_0-1)}$ within
statistical error for $K_0\le6$ and consistent with it beyond; $\Pr[\mathrm{Flip}=\emptyset]=0.125,0.079,0.041\approx4/n$;
the ladder pairs are flippable with frequency $0.5002$ ($62082$ pairs at $n=50$); and
$\Pr[B]-\Pr[A]=-0.022\pm0.011$, $+0.008\pm0.012$, $-0.030\pm0.018$.  On completions of a
fixed rectangle (Remark~\ref{rem:fixrect}; $7\cdot10^5$ marked instances at $n=30$)
$\Pr[\mathrm{Flip}=\emptyset\mid K_0-1\ge K]=2^{-K}$ within noise to $K=12$--$15$, in
every tested type.
\end{remark}

\subsection{Offset ladders and short arcs}\label{sec:offsets}


The ladder is one of a family.  For an integer $c$ (an \emph{offset}) put
\[
K_c=\begin{cases}\min(|C_1|-c,\ |C_2|)&\text{on }A,\\
\min(d_{21}-c,\ d_{12})&\text{on }B,\end{cases}
\qquad
D_c=\bigl\{(x_{c+i},\,y_i)=(\pi^{-(c+i)}(1),\pi^{-i}(2)):\ 1\le i<K_c\bigr\},
\]
so $D_0$ is the ladder and for $c>0$ the $x$-row is $c$ steps deeper
than the $y$-row ($x$-depths $c+1,\dots,c+K_c-1$ stay inside
$C_1\setminus\{1\}$, resp.\ the arc from $2$ to $1$; $y$-depths
$1,\dots,K_c-1$ inside $C_2\setminus\{2\}$, resp.\ the arc from $1$ to
$2$).  We take $c\ge0$; the families with the $y$-row deeper are
obtained by exchanging the roles of rows $1$ and $2$ (and of
$C_1,C_2$, $d_{12},d_{21}$) throughout.

\begin{theorem}[offset ladders]\label{thm:ladderoffset}
For every offset $c\ge0$, on the domain $\mathcal D_c=\{K_c\ge2\}$
(where $D_c\ne\emptyset$) the $j'$-trades at the pairs of $D_c$ are commuting
involutions that preserve $K_c$, the unordered pairs of $D_c$ and their
flippabilities, and toggle the bit iff the traded pair is flippable.
Hence, with $\mathrm{Flip}_c$ the set of flippable pairs of $D_c$,
\[
\Pr[B,\mathcal D_c]-\Pr[A,\mathcal D_c]
=\Pr[B,\mathcal D_c,\mathrm{Flip}_c=\emptyset]-\Pr[A,\mathcal D_c,\mathrm{Flip}_c=\emptyset].
\]
\end{theorem}

\begin{proof}
Identical to Theorem~\ref{thm:ladder}; only the bookkeeping changes.
Let $\pi'=(x_{c+i}\,y_i)\circ\pi$.  The preimage chain of $1$ under
$\pi'$ agrees with that of $\pi$ up to depth $c+i$ and then continues
along the chain of $2$: $\pi'^{-(c+i+t)}(1)=y_{i+t}$; symmetrically
$\pi'^{-(i+t)}(2)=x_{c+i+t}$ ($t\ge1$).  So the pair of $D_c$ at
$y$-depth $i'>i$ is $\{y_{i'},x_{c+i'}\}$, unchanged as a set, and the
pairs at depth $i'<i$ are untouched: $D_c$ is preserved.  (The same
computation shows that trading a pair at depths $(k,k')$ sends a
deeper pair at depths $(K,K')$ to depths $(K'+k-k',\,K-k+k')$, whose
offset is $2(k-k')-(K-K')$; it is preserved iff $K-K'=k-k'$, so only
families of constant offset, i.e.\ the $D_c$, are preserved.)  On $A$
the trade merges $C_1,C_2$ into a cycle with
$d'_{12}=|C_1|-(c+i)+i=|C_1|-c$ and $d'_{21}=|C_2|+c$, so
$K'_c=\min(d'_{21}-c,d'_{12})=\min(|C_2|,|C_1|-c)=K_c$; on $B$ the pair
separates rows $1,2$ ($c+i<d_{21}$ puts $x_{c+i}$ on the arc from $2$
to $1$, $i<d_{12}$ puts $y_i$ on the arc from $1$ to $2$) and gives
$|C'_1|=d_{12}+c$, $|C'_2|=d_{21}-c$, again
$K'_c=\min(d_{12},d_{21}-c)=K_c$.  No pair of $D_c$ is frame-adjacent:
$y_i=\pi(x_{c+i})$ would need $d_{21}=c-1$ (or $\pi(1)=2$ when $c=0$)
and $x_{c+i}=\pi(y_i)$ would need $d_{21}=c+1$, both excluded on
$\mathcal D_c$.  Commutation,
preservation of flippability and the orbit count are as before.
\end{proof}

Exact enumeration at $n=7$ ($\lambda=(7)$, $\{\alpha,\beta\}=\{2,5\}$)
confirms the identities for $c=0,1,2,3$ to all digits
($\Pr[B,\mathcal D_c]-\Pr[A,\mathcal D_c]=-0.005848,\ +0.001462,\
+0.003509,\ -0.001316$, equal to the right-hand sides).

\paragraph{Using the offsets to price short arcs.}  Fix $\ell\ge2$ and
consider $B$-instances with $d_{12}=\ell$.  For $0\le c\le n/2-\ell$
and $d_{21}\ge c+\ell$ one has $K_c=\ell$, so $D_c$ has $\ell-1$ pairs,
and the offset-$c$ trade at a flippable pair of $D_c$ lands in $A$
with $|C_1|=\ell+c\le n/2$ and $|C_2|=d_{21}-c\ge\ell$.  Within a
class of $\mathcal D_c$ all $B$-members have the same $(d_{12},d_{21})$
and all $A$-members the same $(|C_1|,|C_2|)=(d_{12}+c,d_{21}-c)$, and
$\mathrm{Flip}_c$ is a class invariant, so Theorem~\ref{thm:ladderoffset}
class by class and Proposition~\ref{prop:trapped} give
\begin{equation}\label{eq:offsetprice}
\Pr\bigl[B,\ d_{12}=\ell,\ d_{21}\ge c+\ell,\ \mathrm{Flip}_c\ne\emptyset\bigr]
=\Pr\bigl[A,\ |C_1|=\ell+c,\ |C_2|\ge\ell,\ \mathrm{Flip}_c\ne\emptyset\bigr]
\le\frac{18}{n-\ell-c-1}\le\frac{36}{n-2}.
\end{equation}
Averaging over $c$: if $\theta(L,P)$ denotes the fraction of offsets
$c\in[0,n/2-\ell]$ whose diagonal $D_c$ contains a flippable pair,
then for every $\theta_0>0$
\begin{equation}\label{eq:offsetavg}
\Pr\bigl[B,\ d_{12}=\ell,\ |C_1|\ge n/2+\ell,\ \theta\ge\theta_0\bigr]\le\frac{36}{(n-2)\,\theta_0},
\end{equation}
since on $|C_1|\ge n/2+\ell$ one has $d_{21}=|C_1|-\ell\ge n/2\ge c+\ell$
for every such $c$, so $\E[\theta\,\mathbf 1_E]$ is the average over
$c$ of the left side of \eqref{eq:offsetprice}, and Markov's
inequality applies.  Together with the splice-in bound
$\Pr[B,d_{12}=\ell,|C_1|<n/2+\ell]\le4/(n-2\ell+1)\le8/n$ for
$\ell\le n/4$ (Proposition~\ref{prop:splice} below), this prices every short arc except on
the event $\{\theta<\theta_0\}$.


\begin{proposition}[short frame cycles through a marked row are rare]
\label{prop:trapped}
For $(L,P)$ uniform on $X$ and every $2\le\ell\le n/2$,
\[
\Pr\bigl[|C_1|=\ell\bigr]\;\le\;\frac{18}{n-\ell-1},
\qquad\text{hence}\qquad
\Pr\bigl[|C_1|\le \ell\bigr]\;\le\;\frac{36\,\ell}{n}\quad(\ell\le n/2-1),
\]
and the same bounds hold for $C_2$.
\end{proposition}

\begin{proof}
Let $E_\ell=\{(L,P)\in X:\ |C_1|=\ell\}$.  (The case $C_1=\{1,2\}$
cannot occur: it would mean $\sigma(j)=j'$ and $\sigma(j')=j$, but
$j'=\sigma^\alpha(j)$ lies in an $m$-cycle with $m\ge4$.)

\emph{Forward moves.}  From $(L,P)\in E_\ell$ choose $y\in C_1\setminus\{1,2\}$
and $x\notin C_1\cup\{2\}$: at least $(\ell-2)(n-\ell-1)$ pairs
($(\ell-1)(n-\ell-1)$ if $2\notin C_1$; for $\ell=2$ this is $n-3$).
If $(x,y)$ is flippable, apply the $j'$-trade at $(x,y)$ (trade rows $x,y$ along
the $\rho_{x,y}$-cycle $Q_{x,y}$ through $j'$); otherwise turn $C_1$ (which
makes $(x,y)$ flippable, as $y\in C_1\not\ni x$) and then trade
$Q_{x,y}$ in the turned square.  In both cases the resulting frame is
$(x\,y)\circ\pi^\pm$ with $\pi^\pm$ equal to $\pi$ or to $\pi$
inverted on $C_1$, so the cycle of row $1$ becomes $C_1\cup C_x$, of
length $\ell+|C_x|>\ell$; rows $1,2$ are untouched by the trade, and
the result is in $X$ (no turn, or $2\in C_1$) or in $J$ (turn with
$2\notin C_1$).

\emph{Backward moves.}  A square $(L',P)\in X\cup J$ is reached from
$(L,P)\in E_\ell$ by the move at $(x,y)$ only if $x,y$ lie in the cycle
$C'_1$ of row $1$ in $L'$, and undoing the trade (the same trade, since
$\rho'_{x,y}=\rho_{x,y}^{-1}$ on $Q_{x,y}$) splits $C'_1$ into the cycles
$\{x,\pi'(x),\dots,\pi'^{-1}(y)\}$ and $\{y,\pi'(y),\dots,\pi'^{-1}(x)\}$ of $(x\,y)\circ\pi'$, of which the one
containing row $1$ must have length exactly $\ell$: given $L'$, that
arc is determined by its position relative to row $1$ ($\ell$
choices), and the arc determines $\{x,y\}$.  With the binary choice
``turn back or not'' this gives at most $2\ell$ preimages per $(L',P)$.

\emph{Counting.}  The map $((L,P),\text{move})\mapsto((L',P),\text{inverse move})$
is injective, so
$|E_\ell|\,(\ell-2)(n-\ell-1)\le2\ell\,(|X|+|J|)\le6\ell\,|X|$, using
$|J|=2|X_A|\le2|X|$ (Theorem~\ref{thm:marked}).  For $\ell\ge3$ this
is the claim; for $\ell=2$ the forward count is $n-3$ and the backward
count $4$, giving $12/(n-3)$.  The bound for $C_2$ is proved by the
same argument with the roles of rows $1$ and $2$ exchanged: a turn of
$C_2\not\ni1$ replaces $\sigma$ by $(\sigma(j)\,\sigma(j'))\circ\sigma$,
which also splits the $m$-cycle into lengths $\alpha,\beta$ and lands
in $J$.
\end{proof}

\begin{proposition}[the splice-in bound]\label{prop:splice}
For $2\le\ell\le n/4$,
$\Pr_X[B,\ d_{12}=\ell,\ |C_1|<n/2+\ell]\le4/(n-2\ell+1)\le8/n$, and the same with
$d_{21}$ in place of $d_{12}$.
\end{proposition}

\begin{proof}
Let $E$ be the event.  \emph{Forward moves.}  From $(L,P)\in E$ choose $y=\pi^i(1)$ with
$1\le i\le\ell-1$ (a row strictly inside the arc from $1$ to $2$; $\ell-1$ choices) and a row
$x\notin C_1$ ($n-|C_1|\ge(n-2\ell+1)/2$ choices), and let $C_x$ be the frame cycle of $x$,
which contains neither $1$ nor $2$.  If $(x,y)$ is not flippable, turn $C_x$: this trades
columns $j,j'$ on the rows of $C_x$ only, so it preserves $X$, inverts $\pi$ on $C_x$
(Lemma~\ref{lem:cycletrade}), fixes row $y$, and replaces $\rho_{x,y}$ by
$\rho_{x,y}\circ(j\,j')$, which makes $(x,y)$ flippable.  Then apply the $j'$-trade at
$(x,y)$, which replaces the frame by $(x\,y)\circ\pi^{\pm}$ ($\pi^\pm=\pi$, or $\pi$ inverted
on $C_x$): the arc from $1$ to $2$ becomes
$1\to\dots\to\pi^{-1}(y)\to x\to\dots\to y\to\dots\to2$ with $C_x$ occupying
positions $i,\dots,i+|C_x|-1$, so $d_{12}$ grows by $|C_x|$, rows $1,2$ are untouched, and the result
is again in $X_B$ with the same mark.  \emph{Backward moves.}  Given the target
$(L',P)\in X_B$ and $\ell$, a preimage is determined by the position $i\in\{1,\dots,\ell-1\}$
of $x$ on the arc of $L'$ from $1$ to $2$ (then $|C_x|=d'_{12}-\ell$ and
$y=\pi'^{\,i+|C_x|}(1)$ are determined), together with the binary choice of whether $C_x$ was turned; undoing is the
same $j'$-trade (the pair stays flippable, $\rho'_{x,y}=\rho_{x,y}^{-1}$ on the traded
cycle) followed by the turn.  So each target has at most $2(\ell-1)$ preimages while each
source has at least $(\ell-1)(n-2\ell+1)/2$ forward moves, whence
$|E|\,(\ell-1)(n-2\ell+1)/2\le2(\ell-1)\,|X|$.
\end{proof}


\section{The witness hypothesis}\label{sec:witness}

Notation as in \S\ref{sec:ladder}: $(L,P)\in X$, $P=\{j,j'\}$, frame $\pi$, ladder pairs
$(x_k,y_k)$, $K_0$, and for $B$-instances the offset diagonals $D_c$ of
\S\ref{sec:offsets}.  The statement below is the reference hypothesis of this paper --- the
least demanding sufficient condition for (L) that we have; the orbit method of
\S\ref{sec:orbit} gives two others.

Say that a pair of rows
$(x,y)$, $x,y\notin\{1,2\}$, carries a \emph{witness} if the
$\rho_{x,y}$-cycle through column $j$ has length at most
$\ell'=\log^5n$ and does not contain $j'$ (then $(x,y)$ is flippable).
A frame-adjacent pair ($y=\pi(x)$, so $\rho_{x,y}(j)=j'$) never
carries one; the families below contain no such pair.
For a pair of rows of a random Latin square the length of the
$j$-cycle is spread like a uniform variable on $\{2,\dots,n\}$
(Remark~\ref{rem:witnessdata} below), so a witness has probability
$\approx\ell'/n$.

\begin{hypothesis}[witness lemma]\label{hyp:witness}
Let $(L,P)$ be uniform on $X$, uniformly in $\lambda,\alpha,\beta$.  Let
$\mathcal G$ be one of the following families of \emph{blocks} of pairs
of rows, determined by the frame: (i) the ladder, each pair a block of
its own; (ii) for a $B$-instance, with $\ell:=d_{12}$, provided
$\ell\le n/4$ and $|C_1|\ge n/2+\ell$, the diagonals $D_c$,
$0\le c\le n/2-\ell$, each a block of $\ell-1$ pairs (one family, not
one per $\ell$).  Let $G_{\mathcal G}$ be the number of blocks
containing a pair that carries a witness, and
$\mu_{\mathcal G}=\sum_{\mathcal B\in\mathcal G}\min\bigl(\tfrac12,\ell'|\mathcal B|/n\bigr)$.
Then, for the ladder (i),
\[
\Pr\bigl[G_{\mathcal G}=0\ \text{ and }\ \mu_{\mathcal G}\ge\log^2n\bigr]=o(1/\log n),
\]
and for the diagonals (ii) there is a constant $\kappa_1>0$ such that
\[
\Pr\bigl[G_{\mathcal G}<\kappa_1\,\mu_{\mathcal G}\ \text{ and }\ \mu_{\mathcal G}\ge\log^2n\bigr]=o(1/\log n).
\]
(This is annealed: an exceptional set of frames of probability
$o(1/\log n)$ is allowed.  For a single pair of non-adjacent rows the
witness probability is $\approx\ell'/n$, Remark~\ref{rem:witnessdata},
so $\kappa_1\mu_{\mathcal G}$ is a constant fraction of the expected
count if blocks are nearly independent given their shared rows.)
\end{hypothesis}

\begin{proposition}[(L) follows from the witness lemma]\label{prop:Lwitness}
Hypothesis~\ref{hyp:witness} implies \textup{(L)}, hence the weak CGW
conjecture.
\end{proposition}

\begin{proof}
Put $\ell_0=n/\log^2n$.  By Theorem~\ref{thm:ladder},
$|\Pr[B]-\Pr[A]|\le\Pr[\mathrm{Flip}=\emptyset]\le
\Pr[K_0\ge\ell_0,\mathrm{Flip}=\emptyset]+\Pr[A,K_0<\ell_0]+\Pr[B,K_0<\ell_0]$.

\emph{Long ladders.}  On $\{K_0\ge\ell_0\}$ take family (i):
$\mu\ge(\ell_0-1)\ell'/n\ge\log^2n$, so $G\ge1$ except on an
event of probability $o(1/\log n)$, and $G\ge1$ means
$\mathrm{Flip}\ne\emptyset$.

\emph{Short cycles.}  $\Pr[A,K_0<\ell_0]\le72\ell_0/n=72/\log^2n$ by
Proposition~\ref{prop:trapped}.

\emph{Short arcs.}  $\Pr[B,K_0<\ell_0]\le\sum_{\ell<\ell_0}\bigl(\Pr[B,d_{12}=\ell]+\Pr[B,d_{21}=\ell]\bigr)$;
we treat $d_{12}=\ell$; the case $d_{21}=\ell$ reduces to it by exchanging
rows $1$ and $2$ (with the symbol relabelling that restores the first
row): this replaces $\sigma$ by $\sigma^{-1}$, keeps the mark $\{j,j'\}$
with $j'=\sigma^{-\beta}(j)$, so maps $X(n,\lambda;\alpha,\beta)$ onto
$X(n,\lambda;\beta,\alpha)$ measure-preservingly, conjugates the frame
by $(1\,2)$, hence exchanges $d_{12}$ with $d_{21}$, preserves the
ladder as a set of unordered pairs and carries the diagonals with the
$x$-row deeper to those with the $y$-row deeper, and leaves every
$\rho_{x,y}$ with $x,y\notin\{1,2\}$ unchanged; so the $d_{21}=\ell$ case
for $(\lambda,\alpha,\beta)$ is the $d_{12}=\ell$ case of the hypothesis
for $(\lambda,\beta,\alpha)$, which ``uniformly in $\lambda,\alpha,\beta$''
covers (for $\alpha=\beta$ the hypothesis is read for the ordered mark).
Fix $\ell$.  The part with $|C_1|<n/2+\ell$ costs $\le8/n$ (Proposition~\ref{prop:splice}).  On
$|C_1|\ge n/2+\ell$ take family (ii), with $N=\lfloor n/2-\ell\rfloor+1\ge n/4$
blocks: $\mu=N\min(\tfrac12,(\ell-1)\ell'/n)\ge\tfrac14\min(n/2,(\ell-1)\ell')\ge\log^2n$,
so except on an event $F_\ell$ of probability $o(1/\log n)$ the
fraction of offsets whose diagonal contains a flippable pair is
$\theta=G/N\ge\theta_0(\ell):=\kappa_1\min(\tfrac12,(\ell-1)\ell'/n)$.  By
\eqref{eq:offsetavg},
\begin{multline*}
\sum_{\ell<\ell_0}\Pr[B,d_{12}=\ell,|C_1|\ge n/2+\ell,\theta\ge\theta_0(\ell)]
\le\sum_{\ell-1\le n/(2\ell')}\frac{36\,n}{(n-2)\,\kappa_1\,(\ell-1)\ell'}
+\sum_{n/(2\ell')<\ell-1<\ell_0}\frac{72}{(n-2)\kappa_1}\\
=O\Bigl(\frac{\log n}{\ell'}+\frac{\ell_0}{n}\Bigr)=O\Bigl(\frac1{\log^2n}\Bigr),
\end{multline*}
and $\sum_{\ell<\ell_0}8/n\le8/\log^2n$.  The exceptional events: the
family (ii) is a function of the frame, and the events $\{d_{12}=\ell\}$
are disjoint, so $\sum_\ell\Pr[B,d_{12}=\ell,|C_1|\ge n/2+\ell,F_\ell]
=\Pr\bigl[\bigcup_\ell\{B,d_{12}=\ell,|C_1|\ge n/2+\ell,F_\ell\}\bigr]$ is
the probability of a single failure event of the hypothesis, which is
$o(1/\log n)$.
Altogether $\Pr[\mathrm{Flip}=\emptyset]=o(1/\log n)$.
\end{proof}

\begin{remark}[what the witness lemma is, and is not]
(i) It is a statement about substructures of polylogarithmic size
(a cycle of length $\le\log^5n$ of a row-pair permutation through a
given column), in the space of completions of two rows and two
columns, for families of pairs selected by the frame.  It asks for a
lower bound on their expected number and a lower-tail
concentration (a second moment with relative variance $o(1/\log n)$
suffices).  For two \emph{fixed} rows, counts of such substructures
are what the switching method of \cite{KS18} controls; what it does
not address is a family of row pairs selected by the square, and that
is where the difficulty sits (\S\ref{sec:remains}).
(ii) It asks nothing about long cycles, same-cycle events, or the
fairness of any global statistic.  (iii) Rigid squares ($\mathbb Z_p$
isotopes and their relatives, where no row pair has a short cycle)
enter only through the failure probability, which must be
$o(1/\log n)$; for them the hypothesis is simply the statement that
they are rare.  (iv) The constants: $\ell'=\log^5n$ and
$\ell_0=n/\log^2n$ are one admissible choice; the proof needs
$\ell_0\ell'/n\ge\log^2n$, $\ell'\ge4\log^2n$, $\ell_0/n=o(1/\log n)$ and
$\log n/\ell'=o(1/\log n)$.  (v) The block form is what the
short-arc argument uses: the number of offsets with a witness, not the
number of witnesses, since the diagonals of different offsets are
dependent through the shared rows $y_i$.  (vi) The hypothesis is
stated only for the two families used; a version for arbitrary
frame-determined families would be false, since frame-adjacent pairs
never carry a witness.
\end{remark}


\begin{remark}[data]\label{rem:witnessdata}
On Jacobson--Matthews samples with a uniformly random mark (\texttt{runs/s13/f1/}), the
length of the $j$-cycle of $\rho_{x_k,y_k}$ for a ladder pair is spread like a uniform
variable on $\{2,\dots,n\}$: $\Pr[\le\ell]=0.021,0.081,0.18,0.38$ for $\ell=2,5,10,20$ at
$n=50$ against $(\ell-1)/(n-1)=0.020,0.082,0.18,0.39$, and $0.010,0.041,0.092,0.19$ at
$n=100$; among instances with $K_0\ge n/4$ a flippable ladder pair with $j$-cycle of length
$\le10$ exists in $93\%$ ($n=50$) and $95\%$ ($n=100$) of the instances, and one of length
$\le20$ in $99.6\%$ and $99.8\%$.  These samples see only typical types $\lambda$; the
hypothesis is required uniformly in $\lambda$.  The fixed-rectangle sampler of
Remark~\ref{rem:fixrect} reaches the rare types: there the witness events along the ladder
are independent Bernoulli variables to three digits (flat hazard,
$\Pr[\text{no witness in }K\text{ pairs}]=(1-w_1)^K$ to $K=40$ at $n=100$) and the witness
counts across the offset diagonals have binomial lower tails, identically for $(n)$ and for
$(4,2^{(n-4)/2})$.  Classes with all cycles of $\rho_{1,2}$ of length $\le3$ carry no
admissible mark and do not enter.
\end{remark}

\section{The orbit method: alternative sufficient conditions}\label{sec:orbit}

Throughout this section $X=X(n,\lambda;\alpha,\beta)$, and ``legal'' means: preserves the
type of rows $1,2$ and the mark, i.e.\ maps $X$ to itself.  The mark is now written
$\{p,p'\}$ (so $p=j$, $p'=j'$ in the notation of \S\S\ref{sec:marked}--\ref{sec:ladder}; the
letters $j,j'$ are freed for other columns), $\pi$ is its frame, and a column pair $(q,c)$ is
\emph{free} if $\{q,c\}\cap\{p,p'\}=\emptyset$.  A pair of rows $(x,y)$ is called
\emph{parallel} if it is flippable ($p,p'$ in different cycles of $\rho_{x,y}$) and
\emph{crossed} otherwise; $Q_p$ and $Q_{p'}$ denote the $\rho_{x,y}$-cycles through $p$ and
$p'$.  A free pair $(q,c)$ is \emph{separated} by $\{p,p'\}$ in $\rho_{x,y}$ if $q$ and $c$
lie on different arcs of a common cycle through $p,p'$ (crossed case) or one in $Q_p$ and the
other in $Q_{p'}\ne Q_p$ (parallel case).  The method bounds
$\Pr_X[\mathrm{Flip}_I=\emptyset]$ by a product of conditionally independent factors; the
two hypotheses it leads to are compared with Hypothesis~\ref{hyp:witness} after (c$'$), and
in (e) the same mechanism, applied to the first ladder pair alone, gives an exact expression
for $\Pr_X[B]-\Pr_X[A]$ and a third alternative hypothesis.

\paragraph{(a) Transposition: which half of the conditioning is the
obstacle.}  Let $L$ be uniform on \emph{all} Latin squares of order
$n$ (no condition on rows $1,2$ in this paragraph), let $x,y$ carry
four distinct symbols on $\{p,p'\}$, and condition on the two full
columns $p,p'$ (the frame $\pi$ of $(p,p')$).  In the transposed square $L^{\mathsf T}$ the
columns $p,p'$ are the distinguished rows, the frame $\pi$ is their
column permutation, and the event ``$p,p'$ lie in different cycles of
$\rho_{x,y}$'' is the event ``$\{x,y\}$ is an $A$-pair for the
distinguished rows''.  Conditioning on the two full rows $p,p'$ of
$L^{\mathsf T}$ is conditioning on their type and on the position of
the pair $\{x,y\}$ in their cycle structure, which is free by symbol
and column relabelling.  Hence, by Lemma~\ref{lem:switchtoggle}
($|J_A|=|J_B|$, which needs no condition on the cycle lengths) and
CGW's Theorem~3.13 (the inequality $\Pr_X[A]\ge\tfrac13$, whose published proof has a gap for every split other than $(2,2)$, Remark~\ref{rem:gap}):
\begin{equation}\label{eq:CPcols}
\Pr\bigl[p,p'\text{ in different cycles of }\rho_{x,y}\bigm|\text{columns }p,p'\bigr]
\;\begin{cases}=\tfrac12&\text{if $x,y$ are apart in the frame},\\ \ge\tfrac13&\text{if together at distances $\ge2$}.\end{cases}
\end{equation}
(The second bound is CGW's $\tfrac32$, affected by the gap; it holds with $\tfrac14$ unconditionally at distance $2$, Remark~\ref{rem:gap}.)  So the statement with the \emph{columns} prescribed is CGW's theorem; it is the
conditioning on rows $1,2$ --- on their type and the mark, which is
what defines $X$ and cannot be removed --- that is the obstacle.
In the same transposed picture the ladder trades $T_k$ of
Theorem~\ref{thm:ladder} at flippable pairs are CGW's flips and
backflips at the column pairs $\{x_k,y_k\}$ of $L^{\mathsf T}$, and
Theorem~\ref{thm:ladder} is exactly the legal half of CGW's
Lemma~3.12 for the distinguished rows $p,p'$: the other half --- the
switch and the cross-switch at $\{x_k,y_k\}$ --- trades the columns
$p,p'$ of $L$ along a frame cycle or frame arc through exactly one of
rows $1,2$ (apart: $C_1\ni1$, $C_1\not\ni2$; together: the arc from
$x_k$ to $y_k$ contains $1$ but not $2$, the other arc $2$ but not
$1$), so it changes the type of rows $1,2$ and leaves $X$.  The switch
is, in the original orientation, the unflip of rows $1,2$ at $(p,p')$
(which maps $X_A\to J_A$); it toggles the crossed/parallel status of
\emph{every} pair $(x,y)$ with $x$ on the side of row $1$ and $y$ on
the side of row $2$, and reverses the orientation of one side, so it
carries the ladder (a diagonal of the side-by-side array of pairs) to
a different family (an anti-diagonal).  The cross-switch likewise maps
$X_B$ to $J_B$ (it splits the $\sigma$-cycle through $p,p'$ while
keeping rows $1,2$ in one frame cycle, CGW Lemma~3.10) and toggles
the interior pairs.  Both are exact but not usable for the ladder bits
in $X$.


\paragraph{(b) Row trades are free in $X$; column turns are
conditionally free.}  For rows $z,w\notin\{1,2\}$ every trade of rows
$z,w$ along a $\rho_{z,w}$-cycle or $\rho_{z,w}$-arc changes rows
$z,w$ only, hence is legal in $X$.  In particular:

\begin{lemma}[exact fairness of column-pair status for generic rows]
\label{lem:rowfair}
Let $x,y\notin\{1,2\}$ and let $(q,c)$ be a column pair lying in two
different cycles of $\rho_{x,y}$.  The CGW switch at $(q,c)$ for the
rows $x,y$ (trade of rows $x,y$ along the one of the two cycles that
contains the smaller minimum symbol) is a legal involution of $X$
that toggles whether $x,y$ lie in one column cycle of $(q,c)$, and
preserves the cycle partition of $\rho_{x,y}$ and the separation
status of every column pair with respect to $\{p,p'\}$.  Hence,
conditionally on everything determined by the other rows and by the
cycle partition of $\rho_{x,y}$,
\[
\Pr_X\bigl[x,y\text{ in different column cycles of }(q,c)\bigr]=\tfrac12 .
\]
\end{lemma}

\begin{proof}
As in Lemma~\ref{lem:switchtoggle} with $(1,2)$ replaced by $(x,y)$:
the trade swaps the entries of rows $x,y$ on the columns of the chosen
cycle $Q\ni q$, $Q\not\ni c$, so $\delta_{q,c}\mapsto\delta_{q,c}\circ(x\,y)$,
which splits or joins the column cycles through $x,y$.  It inverts $Q$
in place, so the symbol sets of the two cycles are unchanged and the
min-symbol rule re-selects $Q$: an involution.  Rows $1,2$ are
untouched, so type and mark are preserved.  Inverting $Q$ exchanges the
two arcs between $p$ and $p'$ when $Q\ni p,p'$, so separation of any
pair by $\{p,p'\}$ is preserved.  (If $Q$ contains $p$ or $p'$ the
frame changes at rows $x,y$; so frame-determined data are part of the
conditioning only when $Q$ avoids $p,p'$.)
\end{proof}

A column turn, by contrast, changes the rows of a column cycle.  Let
$(q,c)$ be a free column pair and $Z$ one of its column cycles.

\begin{lemma}[legal column turns]\label{lem:legalturn}
The turn of $Z$ (swapping the entries of columns $q,c$ in the rows of
$Z$) is a measure-preserving involution of $X$ if and only if
$|Z\cap\{1,2\}|\in\{0,2\}$; when $|Z\cap\{1,2\}|=2$ it conjugates
$\sigma$ by $(q\,c)$ and keeps the type and the mark.  It preserves the
columns $p,p'$ (hence the frame and every frame-determined family of
row pairs) and the column-cycle partition of $(q,c)$ and of every
column pair disjoint from $\{q,c\}$ (in particular of every pair of
a matching containing $(q,c)$).
For a row pair $(x,y)$ with $|Z\cap\{x,y\}|=1$ it replaces
$\rho_{x,y}$ by $\rho_{x,y}\circ(q\,c)$ (if $x\in Z$) or
$(q\,c)\circ\rho_{x,y}$ (if $y\in Z$), which toggles the
crossed/parallel status of $(x,y)$ exactly when $(q,c)$ is separated
by $\{p,p'\}$ in $\rho_{x,y}$ (different arcs when crossed, different
cycles $Q_p,Q_{p'}$ when parallel), and leaves $(q,c)$ separated
afterwards; with $|Z\cap\{x,y\}|\in\{0,2\}$ the status is unchanged.
Turns of cycles of column pairs from a fixed matching $\mathcal M$ of
the free columns pairwise commute.
\end{lemma}

\begin{proof}
A column cycle is a Latin trade (Lemma~\ref{lem:cycletrade}), so the
turn is a bijection of Latin squares and an involution, and it
preserves the cycle partition of $(q,c)$; column pairs disjoint from
$\{q,c\}$ are untouched, and so are columns $p,p'$.  Rows $1,2$ change
iff $Z$ meets them; if $Z\ni1,2$ both rows swap their entries at
$q,c$, so $\sigma'=(q\,c)\sigma(q\,c)$, whose type is that of $\sigma$
and in which $p,p'$ (fixed by $(q\,c)$) keep their distance.  If $Z$
meets exactly one of rows $1,2$, $\sigma'=\sigma\circ(q\,c)$ or
$(q\,c)\circ\sigma$, which changes the type.  The effect on $\rho_{x,y}$: if $x\in Z$ and $y\notin Z$ the entries of row $x$ at $q,c$ are exchanged, so $\rho'_{x,y}=\rho_{x,y}\circ(q\,c)$, and symmetrically; multiplying by a transposition of two points on one cycle splits it at those points, so $p,p'$ are separated iff $q,c$ lie on different arcs between them, and multiplying by a transposition of two points on different cycles merges the cycles, so $p,p'$ are joined iff $q,c$ lie one in each of their cycles; in both cases $q,c$ lie one on each side afterwards.  Turns of cycles of
disjoint column pairs act on disjoint sets of cells; two cycles of the
same pair have disjoint row sets.
\end{proof}


\paragraph{(c) The product bound.}  Fix a matching $\mathcal M$ of the
free columns, independent of the square, and a frame-determined family
$\mathcal G$ of row pairs $(x,y)$, $x,y\notin\{1,2\}$, whose
$x$-rows are pairwise distinct and disjoint from its $y$-rows (a block
of the ladder, where the pairs are even disjoint; or a union of
diagonals $D_c$ with distinct $x$-rows, which may share $y$-rows).
For $(x,y)\in\mathcal G$ call a column pair $(q,c)\in\mathcal M$ a
\emph{clean coordinate} of $(x,y)$ if the column cycle $Z$ of $(q,c)$
through $x$ has length at most $\ell_1$ (a fixed constant), satisfies
$|Z\cap\{1,2\}|\in\{0,2\}$, and contains no row of $\mathcal G$
other than $x$ (no other $x$-row and no $y$-row, in particular not
$y$).  Let $\mathcal T_{x,y}$ be the set of clean coordinates of
$(x,y)$ and $\eta_{x,y}\in\{0,1\}^{\mathcal T_{x,y}}$.  By
Lemma~\ref{lem:legalturn} the turns of all clean coordinates of all
pairs of $\mathcal G$ commute, generate a group acting freely on $X$,
preserve $\mathcal G$ and the sets $\mathcal T_{x,y}$, and
\[
\rho_{x,y}(\eta)\;=\;\rho_{x,y}(0)\circ\!\!\prod_{(q,c)\in\mathcal T_{x,y}:\ \eta_{x,y}(q,c)=1}\!\!(q\,c),
\]
the bit $b_{x,y}(\eta)=[\,(x,y)\text{ parallel}\,]$ depending on
$\eta_{x,y}$ only.  Let the \emph{orbit weight} of $(x,y)$ be
\[
\bar y_{x,y}\;=\;\E_{\eta_{x,y}}\Bigl[\frac{\#\{(q,c)\in\mathcal T_{x,y}:\ (q,c)\text{ separated by }\{p,p'\}\text{ in }\rho_{x,y}(\eta)\}}{|\mathcal T_{x,y}|}\Bigr]
\]
($0$ if $\mathcal T_{x,y}=\emptyset$), a function of the orbit, and
$Y_{\mathcal G}=\sum_{(x,y)\in\mathcal G}\bar y_{x,y}$.

\begin{proposition}[product bound]\label{prop:orbit}
Conditionally on the orbit, the bits $b_{x,y}$, $(x,y)\in\mathcal G$,
are independent, and
\[
\Pr\bigl[(x,y)\text{ crossed}\bigm|\text{orbit}\bigr]\le1-\tfrac12\bar y_{x,y},
\qquad
\Pr\bigl[\text{no pair of }\mathcal G\text{ is flippable}\bigm|\text{orbit}\bigr]
\le\prod_{(x,y)\in\mathcal G}\bigl(1-\tfrac12\bar y_{x,y}\bigr)\le e^{-Y_{\mathcal G}/2}.
\]
Consequently $\Pr_X[\mathrm{Flip}_{\mathcal G}=\emptyset]\le\E_X[e^{-Y_{\mathcal G}/2}]$.
\end{proposition}

\begin{proof}
The action is free: distinct column pairs act on disjoint cells, and
two clean coordinates of different pairs at the same $(q,c)$ are the
cycles through two distinct $x$-rows, each avoiding the other, hence
different cycles; the sets $\mathcal T_{x,y}$ are orbit-invariant by
Lemma~\ref{lem:legalturn} (row sets of the cycles, their meeting
$\{1,2\}$, and the frame-determined family are preserved).  The
uniform measure on $X$ restricted to an orbit is therefore uniform on
the orbit, which is the product of the cubes
$\{0,1\}^{\mathcal T_{x,y}}$.  A clean coordinate of $(x,y)$ contains
$x$ and not $y$, and a clean coordinate of another pair contains
neither $x$ nor $y$ (it avoids all $x$-rows but its own and all
$y$-rows), so $b_{x,y}$ depends on $\eta_{x,y}$ alone, whence the
independence.  Fix $(x,y)$ and
$Z=(q,c)\in\mathcal T_{x,y}$.  Whether $Z$ is separated in
$\rho_{x,y}(\eta)$ does not depend on $\eta_Z$
(Lemma~\ref{lem:legalturn}: the turn of $Z$ leaves $Z$ separated), and
when it is, the two values of $\eta_Z$ give opposite bits.  Hence
$\Pr[b_{x,y}=0]\le1-\tfrac12\Pr_\eta[Z\text{ separated}]$ for every
$Z$, and averaging the right side over $Z\in\mathcal T_{x,y}$ gives
the first bound; the second is the product of the first over the
independent bits.
\end{proof}



\begin{hypothesis}[orbit genericity; fixed matching]\label{hyp:orbitgen}
Fix the matching $\mathcal M$ of the free columns in consecutive pairs (in the natural
order of the columns) and a constant $\ell_1$ as in (c),
and for a family $\mathcal G$ and $D\subseteq\mathcal G$ write
$Y_D=\sum_{(x,y)\in D}\bar y_{x,y}$ (orbit weights computed with cleanliness with respect to
all rows of $\mathcal G$).  There is a constant $\kappa_0>0$ such that, with
$\ell_0=n/\log^2n$ and uniformly in $\lambda,\alpha,\beta$: (i) for the sub-ladder $I$ of
the first $\min(K_0-1,\lceil n/8\rceil)$ ladder pairs,
$\E_X\bigl[e^{-Y_I/2};\ K_0\ge\ell_0\bigr]=o(1/\log n)$; (ii) for a $B$-instance, with
$\ell:=d_{12}<\ell_0$ and $|C_1|\ge n/2+\ell$, $S=\{2\ell t:\ 0\le t\le(n/2-\ell)/(2\ell)\}$
and $\mathcal G=\bigcup_{c\in S}D_c$: the event that fewer than $\kappa_0|S|$ of the
diagonals $D_c$, $c\in S$, have $Y_{D_c}\ge\kappa_0(\ell-1)$ has probability $o(1/\log n)$
(one event, not one per $\ell$); and symmetrically for $d_{21}$.  (A fixed small fraction,
not a half: with a fixed matching the orbit weight of a single pair vanishes with
probability $\approx0.7$, \S(d), so for $\ell=2$ only a minority of the one-pair diagonals can
qualify.)
\end{hypothesis}
\noindent Since $\bar y_{x,y}$ is typically of constant size (\S(d): $0.16$--$0.19$ on
average with the fixed matching), (i) asks that the sum of $\approx n/8$ such weights not collapse; its content is a
\emph{genericity constant} --- that a clean coordinate is separated by the mark with
probability bounded below --- which none of the exact tools of this paper provides
(\S\ref{sec:remains}).

\paragraph{(c$'$) Adaptive coordinates: the toggling is automatic.}
Proposition~\ref{prop:orbit} needs the clean coordinates of a ladder pair to be
\emph{separated} by $\{p,p'\}$ with positive probability, a genericity input we cannot prove
(\S\ref{sec:remains}).  It disappears if the coordinates are chosen along $\rho_{x,y}$ from
the mark.  For a pair $(x,y)$ of the family and $t\ge1$ let
\[
(q^t,c^t):=\bigl(\rho_{x,y}^{\,t}(p),\ \rho_{x,y}^{\,t}(p')\bigr),
\]
and let $\mu_{x,y}:=\min(d(p\to p'),d(p'\to p))$ if $(x,y)$ is crossed and
$\min(|Q_p|,|Q_{p'}|)$ if parallel.  For $t<\mu_{x,y}$ the pair $(q^t,c^t)$ is free
($q^t\ne c^t$, both $\notin\{p,p'\}$) and separated by $\{p,p'\}$ in $\rho_{x,y}$ by
construction: on different arcs when crossed, one in $Q_p$ and one in $Q_{p'}$ when parallel.
So the turn of its column cycle $Z$ through $x$ (with $y\notin Z$) toggles the status of
$(x,y)$ whenever it is legal (Lemma~\ref{lem:legalturn}).  What makes these coordinates
usable in the orbit method is that they are invariant under their own turn.  (We keep (c)
for comparison: its orbit weight is what the data of (d) measure, and it is where the
genericity constant of \S\ref{sec:remains} lives.)


\begin{lemma}[the adaptive coordinate is orbit-invariant]\label{lem:adaptinv}
Let $t<\mu_{x,y}$, let $Z$ be the column cycle of $(q^t,c^t)$ through
$x$ with $y\notin Z$, and let $L'$ be the turn of $Z$.  Then
$\rho'_{x,y}=\rho_{x,y}\circ(q^t\,c^t)$, the status of $(x,y)$ toggles,
$\mu'_{x,y}=\mu_{x,y}$, and $(\rho'^{\,s}(p),\rho'^{\,s}(p'))=(q^s,c^s)$
for every $s\le t$; in particular the $t$-th candidate of $L'$ is
$(q^t,c^t)$ again, and its column cycle through $x$ is $Z$ (same row
set).
\end{lemma}
\begin{proof}
Write $\rho=\rho_{x,y}$, $q=q^t$, $c=c^t$.  The turn exchanges the
entries of $x$ at $q,c$ (and those of the other rows of $Z$, which do
not affect $\rho$), so $\rho'=\rho\circ(q\,c)$: $\rho'(q)=\rho(c)$,
$\rho'(c)=\rho(q)$, $\rho'=\rho$ elsewhere.  Crossed case, with
$p=u_0\to u_1\to\dots\to u_a=p'\to v_1\to\dots\to v_b=p$ and
$q=u_t$, $c=v_t$ ($t<\min(a,b)$): the cycles of $\rho'$ through $p$
and $p'$ are $u_0\to\dots\to u_t\to v_{t+1}\to\dots\to v_b=p$ and
$u_{t+1}\to\dots\to u_a=p'\to v_1\to\dots\to v_t\to u_{t+1}$, of
lengths $b$ and $a$, so $(x,y)$ is now parallel with
$\mu'=\min(a,b)=\mu$, and the paths $p\to u_s$, $p'\to v_s$ for
$s\le t$ are untouched.  Parallel case, $Q_p=(p=u_0,u_1,\dots,u_{a-1})$,
$Q_{p'}=(p'=v_0,\dots,v_{b-1})$, $q=u_t$, $c=v_t$ ($t<\min(a,b)$): the
single cycle of $\rho'$ is
$u_0\to\dots\to u_t\to v_{t+1}\to\dots\to v_{b-1}\to v_0=p'\to v_1\to\dots\to v_t\to u_{t+1}\to\dots\to u_{a-1}\to u_0$,
with $d'(p\to p')=b$, $d'(p'\to p)=a$, so crossed with $\mu'=\mu$, and
again the paths up to step $t$ are untouched.  The row set of the
turned cycle is preserved (Lemma~\ref{lem:cycletrade}).
\end{proof}

\begin{proposition}[orbit bound with adaptive coordinates]\label{prop:adaptive}
Let $(L,P)\in X$, $I$ a frame-determined family of pairs $(x_k,y_k)$
as in \S(c) (distinct $x$-rows, disjoint from the $y$-rows, none in
$\{1,2\}$; the ladder, or a union of offset diagonals), and
$T,\ell_1\ge1$.  Call the candidate $(k,t)$ \emph{good} if
$t\le\min(T,\mu_k-1)$ and the column cycle $Z_k^t$ of $(q_k^t,c_k^t)$
through $x_k$ has length $\le\ell_1$, meets $\{1,2\}$ in $0$ or $2$
rows and meets the rows of $I$ only in $x_k$.  Let $t_0(k)$ be
the least good $t$ (if any), $C_k=\{q_k^t,c_k^t:\ t\le t_0(k)\}$
(all candidates $t\le\min(T,\mu_k-1)$ if there is none), and
\[
U=\bigl\{k\in I:\ t_0(k)\text{ exists and }\{q_k^{t_0(k)},c_k^{t_0(k)}\}\cap C_l=\emptyset\ \text{for all }l\in I\setminus\{k\}\bigr\}.
\]
Then the turns $T_k$ of $Z_k^{t_0(k)}$, $k\in U$, are legal, commute,
generate a free action of $\{0,1\}^U$ on $X$ that preserves the frame,
the family, every $\mu_l$, $t_0(l)$, $C_l$ ($l\in I$) and $U$; $T_k$
toggles the status of $(x_k,y_k)$ and fixes that of every other pair.
Hence
\[
\Pr_X[\mathrm{Flip}_I=\emptyset]\;\le\;\E_X\bigl[2^{-|U|}\bigr].
\]
\end{proposition}
\begin{proof}
$T_k$ changes only columns $q_k,c_k:=q_k^{t_0(k)},c_k^{t_0(k)}$ in the
rows of $Z_k$ (lengths of cycles are preserved with their row sets, so
the length condition is invariant along with the rest).  Legality: $Z_k$ meets $\{1,2\}$ in $0$ or $2$ rows and
$\{q_k,c_k\}\cap\{p,p'\}=\emptyset$, so the type of rows $1,2$ and the
mark are preserved (Lemma~\ref{lem:legalturn}); columns $p,p'$ are
untouched, so the frame and the ladder are preserved.  For $k\ne l$ in
$U$ the column pairs $\{q_k,c_k\}$, $\{q_l,c_l\}\subseteq C_l$ are
disjoint, so the turns act on disjoint cells and commute, and the
action is free.  Invariance of the data of a pair $l\ne k$ under $T_k$:
rows $x_l,y_l\notin Z_k$, so $\rho_l$, its status, $\mu_l$ and all its
candidates are unchanged; the goodness of candidate $(l,t)$ depends on
the column cycle of $(q_l^t,c_l^t)$ through $x_l$, i.e.\ on the
contents of columns $q_l^t,c_l^t$, which are not in $\{q_k,c_k\}$ for
$t\le t_0(l)$ (resp.\ for all candidates if $t_0(l)$ does not exist) by
the defining condition of $U$; hence $t_0(l)$ and $C_l$ are unchanged.
Invariance of the data of $k$ itself: by Lemma~\ref{lem:adaptinv},
$\mu_k$ and the candidates $t\le t_0(k)$ are unchanged; the candidates
$t<t_0(k)$ have columns different from $q_k,c_k$ (distinct points of
the two paths), so their cycles and non-goodness are unchanged, and
$Z_k$ itself keeps its row set, so $t_0(k)$ is still good; $C_k$ is
unchanged.  If $Z_k\ni1,2$, rows $1,2$ change at $q_k,c_k$, which
affects none of the above.  Hence $U$ is invariant, every orbit of the
cube is a cube of size $2^{|U|}$ on which the status of pair $k\in U$
is $\mathrm{bit}_k(L)\oplus\varepsilon_k$, and $\mathrm{Flip}_I=\emptyset$
forces one value of $(\varepsilon_k)_{k\in U}$; $X$ is uniform.
\end{proof}


\begin{hypothesis}[short clean cycles at the ladder rows]\label{hyp:adaptive}
There are $T,\ell_1=O(\log^2n)$ and a constant $\kappa_2>0$ such that, uniformly in
$\lambda,\alpha,\beta$, with $\ell_0=n/(16\log^2n)$, $M=\lfloor n/\log^4n\rfloor$ and $U$ as
in Proposition~\ref{prop:adaptive}:
\begin{itemize}[leftmargin=2em]
\item[(i)] for the sub-ladder $I$ of the first $\min(K_0-1,\lceil\ell_0\rceil)$ ladder
pairs, $\E_X\bigl[2^{-|U|};\ K_0\ge\ell_0\bigr]=o(1/\log n)$;
\item[(ii)] for a $B$-instance, with $\ell:=d_{12}<\ell_0$ and $|C_1|\ge n/2+\ell$: let
$S=\{2\ell t:\ 0\le t\le(n/2-\ell)/(2\ell)\}$ (offsets at gaps $2\ell$, so that the $x$-rows
of the diagonals $D_c$, $c\in S$, are pairwise distinct; $|S|\ge n/(8\ell)\ge2\log^2n$),
let $S'$ be the first $m_\ell:=\min\bigl(|S|,\lceil n/(\log^2n\cdot\min(\ell-1,M))\rceil\bigr)$
offsets of $S$, let $D_c^{(M)}$ be the first $\min(\ell-1,M)$ pairs of $D_c$, and take $U$
for the family $\mathcal F_\ell=\bigcup_{c\in S'}D_c^{(M)}$ (which has at most
$2n/\log^2n$ pairs).  Then the event that fewer than $\kappa_2\log^2n$ of the diagonals
$D_c^{(M)}$, $c\in S'$, meet $U$ has probability $o(1/\log n)$ (one event, not one per
$\ell$); and symmetrically for $d_{21}$.
\end{itemize}
\end{hypothesis}
\noindent The scaling in (ii) is forced: a diagonal of $\ell-1$ pairs meets $U$ with
probability $\approx\min(\ell-1,M)\,T\ell_1/n$, so for short arcs only a small fraction of
the diagonals can be expected to meet $U$, and the family must be kept to
$O(n/\log^2n)$ pairs for the cleanliness and collision conditions of
Proposition~\ref{prop:adaptive} to have constant cost; $|S'|\min(\ell-1,M)\approx n/\log^2n$
makes the expected number of diagonals meeting $U$ of order $\log^2n$ for every $\ell$.

\begin{proposition}\label{prop:Lorbit}
Each of Hypotheses~\ref{hyp:orbitgen} and~\ref{hyp:adaptive} implies \textup{(L)}, hence
the weak CGW conjecture.
\end{proposition}

\begin{proof}
By Theorem~\ref{thm:ladder},
$|\Pr[B]-\Pr[A]|\le\Pr[\mathrm{Flip}=\emptyset]\le\Pr[K_0\ge\ell_0,\mathrm{Flip}=\emptyset]+\Pr[A,K_0<\ell_0]+\Pr[B,K_0<\ell_0]$,
with $\ell_0$ as in the hypothesis used.  \emph{Long ladders.}  $K_0$ is a function of the
frame, hence orbit-invariant, and by Proposition~\ref{prop:orbit} resp.\
Proposition~\ref{prop:adaptive} $\Pr[\mathrm{Flip}_I=\emptyset\mid\text{orbit}]\le e^{-Y_I/2}$
resp.\ $2^{-|U|}$, so $\Pr[K_0\ge\ell_0,\mathrm{Flip}=\emptyset]=o(1/\log n)$ by (i).
\emph{Short ladders on $A$.}  $\Pr[A,K_0<\ell_0]\le72\ell_0/n$ by Proposition~\ref{prop:trapped}.
\emph{Short arcs on $B$.}  Fix $\ell<\ell_0$ and consider $B$-instances with $d_{12}=\ell$;
the case $d_{21}=\ell$ reduces to it by the involution exchanging rows $1,2$, as in
Proposition~\ref{prop:Lwitness}.  The part $|C_1|<n/2+\ell$ costs at most $8/n$ by
Proposition~\ref{prop:splice}.  On $|C_1|\ge n/2+\ell$ the family of (ii) is as required by
the propositions: its $x$-rows lie in the disjoint depth intervals $[c+1,c+\ell-1]$, $c\in S$,
on the frame arc from $2$ to $1$, and its $y$-rows $\pi^{-i}(2)$ on the arc from $1$ to $2$.
Given the orbit the statuses of all pairs are independent, and the diagonals'
flippabilities are independent (they are disjoint sets of pairs).  Under
Hypothesis~\ref{hyp:orbitgen}, a diagonal with $Y_{D_c}\ge\kappa_0(\ell-1)$ is flippable
with conditional probability $\ge1-e^{-\kappa_0(\ell-1)/2}\ge4\theta_0$,
$\theta_0:=\tfrac14(1-e^{-\kappa_0/2})$; off the exceptional event of (ii) at least
$\kappa_0|S|$ diagonals qualify, so the number of flippable offsets has conditional mean
$\ge4\theta_0\kappa_0|S|$ and the fraction $\theta_S$ is $\ge2\theta_0\kappa_0=:\theta_0'$
except with conditional probability $\le e^{-\theta_0\kappa_0|S|/2}\le e^{-\theta_0\kappa_0\log^2n/16}$
(as $|S|\ge n/(8\ell_0)=\log^2n/8$).
Under Hypothesis~\ref{hyp:adaptive}, a diagonal $D_c^{(M)}$ meeting $U$ is flippable with
conditional probability $\ge\tfrac12$ (the status of a pair $k\in U$ is
$\mathrm{bit}_k\oplus\varepsilon_k$ with $\varepsilon$ uniform on the orbit); off the
exceptional event of (ii) at least $\kappa_2\log^2n$ diagonals of $S'$ meet $U$, so the
number of flippable offsets in $S'$ has conditional mean $\mu\ge\tfrac12\kappa_2\log^2n$ and
is $\ge\mu/2$ except with conditional probability $\le e^{-\mu/8}\le e^{-\kappa_2\log^2n/16}$;
hence the fraction $\theta_{S'}$ is $\ge\theta_0(\ell):=\kappa_2\log^2n/(4m_\ell)$.  In both
cases the conditional failure probabilities are summable over $\ell<\ell_0$ to $o(1/n)$.
The derivation of \eqref{eq:offsetavg} from \eqref{eq:offsetprice} applies to the average
over any set of offsets determined by $\ell=d_{12}$ alone, such as $S$ and $S'$, so
\[
\sum_{\ell<\ell_0}\Pr\bigl[B,d_{12}=\ell,|C_1|\ge n/2+\ell,\ \theta_{S}\ge\theta_0(\ell)\bigr]
\le\sum_{\ell<\ell_0}\frac{36}{(n-2)\,\theta_0(\ell)}
\]
(with $S'$ and $\theta_{S'}$ in the adaptive case), which is
$36\ell_0/((n-2)\theta_0')=O(1/\log^2n)$ for Hypothesis~\ref{hyp:orbitgen}, and
for Hypothesis~\ref{hyp:adaptive}, with $m_\ell\le n/(\log^2n\min(\ell-1,M))+1$,
\begin{multline*}
\sum_{\ell<\ell_0}\frac{144\,m_\ell}{(n-2)\kappa_2\log^2n}
\le\frac{144}{\kappa_2\log^4n}\Bigl(\sum_{\ell-1\le M}\frac1{\ell-1}+\sum_{M<\ell-1<\ell_0}\frac1M\Bigr)+\frac{144\,\ell_0}{(n-2)\kappa_2\log^2n}\\
=O\Bigl(\frac{\log n}{\log^4n}+\frac{\ell_0}{M\log^4n}+\frac1{\log^4n}\Bigr)=O\Bigl(\frac1{\log^2n}\Bigr).
\end{multline*}
The exceptional events of (ii) over all $\ell$ form a single event of probability
$o(1/\log n)$ ($\ell=d_{12}$ is a frame statistic).  Altogether
$\Pr[\mathrm{Flip}=\emptyset]=o(1/\log n)$.
\end{proof}

\paragraph{What the hypothesis says, and what it costs.}  A ladder pair enters $U$ when one
of its $T$ candidate column pairs has a short column cycle through $x_k$ that is clean ---
a configuration of $O(\log^2n)$ cells ($4t$ cells of rows $x_k,y_k$ to locate the
candidate, and the cycle itself).  Heuristically a candidate's column cycle has length
$\le\ell_1$ with probability $\approx\ell_1/n$ (the data of \S(d) show the lengths spread
like a uniform variable on $\{2,\dots,n\}$), is clean with probability
$\approx e^{-2s\ell_1/n}$ for a family of $s$ pairs, and the chosen pair avoids the
$\approx2Ts$ candidate columns of the other pairs with probability $\approx(1-2Ts/n)^2$, so
\[
\E|U|\;\approx\;s\cdot\frac{T\ell_1}{n}\,e^{-2s\ell_1/n}\Bigl(1-\frac{2Ts}n\Bigr)^2\;\asymp\;\log^2n
\qquad\text{for }s=\frac n{16\log^2n},\ T=\frac{\log^2n}4,\ \ell_1=\frac{\log^2n}2,
\]
against the $C\log\log n$ ($C>1/\log2$) that (i) needs; the constant is small
($\approx1/145$), so the heuristic excess over the threshold is a large-$n$ statement.  Compared with the fixed matching
of (c), which offers $\approx n/2$ coordinates per pair of which a constant number are short
and clean and a fraction $\approx0.19$ of those separated, the adaptive rule offers only
$T=O(\log^2n)$ coordinates per pair, all separated: the genericity constant is gone, and in
exchange membership in $U$ is a rare event per pair, so (i) is a lower-tail statement about
rare local events at rows selected by the square --- of the same shape as a statement that
some ladder pair has a short $\rho_{x_k,y_k}$-cycle through $p$ avoiding $p'$ (which would
make it flippable outright), now with a column cycle of a free pair through $x_k$ in its
place.  There is a second regime:
with $s\approx T\approx\ell_1\approx\sqrt n$ a pair enters $U$ with probability bounded
below and $\E|U|\asymp\sqrt n$, so (i) becomes a second-moment statement about $\sqrt n$
events of constant probability, at the price of the arcs of $\rho_k$ at $p,p'$ exceeding
$\sqrt n$ for a constant fraction of the pairs.  In both regimes the lengths of the arcs of $\rho_k$ at $p,p'$ --- which the
fixed-matching version needed to be macroscopic for a positive orbit weight --- enter only
through the number of candidates $\min(T,\mu_k-1)$.  The collision rule in $U$ (the chosen pair must avoid the candidate
columns of \emph{all} other pairs) is what makes $U$ orbit-invariant --- a turn can shrink
the candidate set of an excluded pair, so the obvious greedy refinement is not invariant ---
and it is crude: it is why the asymptotic regime $2Ts/n\le\tfrac12$ is out of reach of
sampling at $n\le100$.

\paragraph{Comparison with the witness lemma.}  Hypothesis~\ref{hyp:witness} asks for
\emph{one} witness on a long ladder (a short $\rho_{x_k,y_k}$-cycle through $p$ avoiding
$p'$, which makes the pair flippable with no further argument), at a per-pair rate
$\approx\ell'/n$.  Hypothesis~\ref{hyp:adaptive} asks for $C\log\log n$ good candidates at
a per-pair rate $\approx T\ell_1/n$, and Hypothesis~\ref{hyp:orbitgen} asks for a
genericity constant.  Neither orbit hypothesis is implied by the witness lemma or implies
it (we see no implication in either direction); the witness lemma is the least demanding
heuristically and is the reference hypothesis of this paper.  What the orbit method contributes is the mechanism ---
exact fairness compounded over commuting involutions into exponential decay --- and, with
the adaptive coordinates, the observation that the mechanism needs only rare local
witnesses of a second kind (column cycles of a free pair through $x_k$) rather than a
constant; at the first ladder pair alone, (e) below, the rarity disappears too, and what is
left is the lower tail of the number of good candidates among $\mu-1$ column pairs.  All of
these hypotheses share the same unproved core, stated in \S\ref{sec:remains}.

\paragraph{(d) Data.}  All scripts and logs are in the repository \cite{repo}, under
\texttt{s13\_*.py} and \texttt{runs/s13/orbit/}.  \emph{Adaptive coordinates.}  On
instances with a uniformly random mark at $n=30,50,100$ (family the first $\min(K_0-1,n/8)$ ladder pairs,
$T=6$, $\ell_1=6,8,10$; $2400$, $3000$, $900$ instances), for a random $k\in U$ the turn of
$Z_k$ was applied and checked to toggle the status of $(x_k,y_k)$, to leave the status,
$\mu_l$, $t_0(l)$, $C_l$ of every other pair and the set $U$ unchanged, and to preserve the
type of rows $1,2$ and the frame ($682$, $697$, $215$ turns, no failure; $444$ and $160$
more with families of $12$ and $30$ pairs); the bound
$\Pr[\mathrm{Flip}_I=\emptyset]\le\E[2^{-|U|}]$ holds in every $K_0$ class.  At these sizes
$30\%$, $23\%$, $16\%$ of the ladder pairs have a good candidate and $13\%$, $5\%$, $3\%$
are in $U$ ($8\%$ and $1.4\%$ with the family of $30$ pairs at $n=100$): $2Ts\ge n$ there,
so the collision rule removes most pairs.  \emph{Column-cycle lengths
and the fixed-matching coordinates.}  The length of the column cycle of a free pair through
a ladder row is spread like a uniform variable on $\{2,\dots,n\}$ (about $1/n$ per length;
the turn is legal for $98\%$ of the $2$-cycles and $64\%$ of the long cycles, and
avoids $y$ for $97\%$ resp.\ $27\%$ of them, at $n=100$).  For the fixed matching of
(c): $\Pr[(q,c)\text{ separated}\mid\text{clean, length }\ell]=0.193$--$0.201$ for
$\ell=2,\dots,8$ against $0.195$ over all column pairs ($n=100$), i.e.\ no detectable
dependence on the cycle length; with $|\mathcal T|\approx42$ coordinates the orbit average
of the parallel bit is $0.478$ from a crossed and $0.523$ from a parallel start ($n=100$);
with the fixed matching of consecutive free columns $|\mathcal T|=1.7$ resp.\ $3.0$ clean
coordinates per pair ($n=50,100$) and $\E[\bar y]=0.159$ resp.\ $0.189$; and the orbit
weights of different ladder pairs are uncorrelated to within sampling error
($\mathrm{Var}(Y_I)/(|I|\mathrm{Var}\,\bar y)=0.99$ and $1.14$ at $n=50,100$, the excess
coming from the variation of $|I|$).  The ladder data of Remark~\ref{rem:ladderdata}
complete the picture: $\Pr[\mathrm{Flip}=\emptyset]\approx4/n$, as if the ladder coins
were fair and independent.


\paragraph{(e) The first ladder pair alone.}  The adaptive coordinates of (c$'$)
specialise to the single pair $(x_1,y_1)=(\pi^{-1}(1),\pi^{-1}(2))$ with none of the side
conditions of Proposition~\ref{prop:adaptive} (there are no other pairs to protect), and
the resulting involution combines with the first ladder trade $T_1$ into a bound on
$\Pr_X[B]-\Pr_X[A]$ by an event selected by that one pair.  Since $\alpha,\beta\ge2$ the
pair always consists of rows outside $\{1,2\}$: $x_1=2$ would mean $p'=\sigma(p)$ and
$y_1=1$ would mean $p=\sigma(p')$ (so $K_0\ge2$ always).  Write $\rho=\rho_{x_1,y_1}$, $F$
for the event that $(x_1,y_1)$ is flippable (parallel), and $\nu$ for the minimum of the
two arc lengths from $p$ to $p'$ and from $p'$ to $p$ (crossed case) or of $|Q_p|,|Q_{p'}|$
(parallel case) --- the quantity written $\mu_{x_1,y_1}$ in (c$'$); we use $\nu$ here
because $\mu$ is the split type elsewhere in the paper.

\begin{lemma}\label{lem:mu2}
$\nu\ge2$ always.
\end{lemma}
\begin{proof}
$\rho(p)=p'$ would mean $L(y_1,p')=L(x_1,p)=L(1,p')$, impossible in a column; $\rho(p')=p$
would mean $L(x_1,p')=L(y_1,p)=L(2,p')$, again impossible in a column; and $\rho$ has no
fixed points since $x_1\ne y_1$.
\end{proof}

For $1\le t<\nu$ the column pair $\{q_t,c_t\}=\{\rho^t(p),\rho^t(p')\}$ (written
$q^t,c^t$ in (c$'$)) is free, consists of two distinct columns, and is separated by the
mark, as observed in (c$'$).  Let $\tau_t=\delta_{q_t,c_t}$ be the permutation of rows
induced by the column pair ($\tau_t(r)$ is the row of $L(r,q_t)$ in column $c_t$) and
$Z_t(r)$ its cycle through row $r$.  Call $t$ \emph{good} if $x_1$ and
$y_1$ lie in different cycles of $\tau_t$ and at least one of $Z_t(x_1)$, $Z_t(y_1)$ meets
$\{1,2\}$ in $0$ or $2$ rows (so that its turn is legal); let
\[
G=\{\text{some }t<\nu\text{ is good}\},
\]
and on $G$ let $t_0$ be the least good $t$ and $Z$ the admissible cycle at $t_0$ (the one
through $x_1$ if it is admissible, otherwise the one through $y_1$).

\begin{proposition}[single-pair involution]\label{prop:firstpair}
Turning $Z$ (exchanging the entries of columns $q_{t_0},c_{t_0}$ in the rows of $Z$) is an
involution of $X\cap G$ that preserves the frame $\pi$ (hence $x_1,y_1$, $A$ and $B$),
toggles $F$, and fixes $t_0$ and $Z$.  Consequently
$\Pr_X[F\cap G]=\Pr_X[F^c\cap G]$ and
$\bigl|\Pr_X[F]-\tfrac12\bigr|\le\tfrac12\Pr_X[G^c]$.
\end{proposition}

\begin{proof}
Columns $p,p'$ are untouched, so $\pi$, $x_1,y_1$, the mark and the rows outside $Z$ are
unchanged.  Rows $1,2$ are both in $Z$ or both outside; if both are in, they become
$L_1\circ(q\,c)$ and $L_2\circ(q\,c)$ (row $1$ is then no longer the identity, which the
unnormalised definition of $X$ allows), so the new row permutation is
$(q\,c)\,\sigma\,(q\,c)$, of the type of $\sigma$, and as $q,c\notin\{p,p'\}$ it still has
$p'$ at distance $\alpha$ from $p$: the turned square is in $X$.  Exactly one of $x_1,y_1$
is in $Z$ (they lie in different $\tau_{t_0}$-cycles), so $\rho'=\rho\circ(q\,c)$ if
$x_1\in Z$ and $\rho'=(q\,c)\circ\rho$ if $y_1\in Z$.  In the first case
Lemma~\ref{lem:adaptinv} applies with $t=t_0$.  In the second,
$(q\,c)\circ\rho=\rho\circ(u\,u')$ with $u=\rho^{t_0-1}(p)$, $u'=\rho^{t_0-1}(p')$
(possibly $t_0=1$ and $\{u,u'\}=\{p,p'\}$), and the computation in the proof of
Lemma~\ref{lem:adaptinv}, read with $t_0-1$ in place of $t$ (for $t_0=1$ the paths are
empty and the computation is immediate), shows that the $\rho$-paths from $p$ to $u$ and
from $p'$ to $u'$ are unchanged, that $\rho'(u)=c$ and $\rho'(u')=q$ (the
ordered pair at step $t_0$ is reversed, the unordered pair is $\{q,c\}$), and that crossed
with arcs $(a,b)$ becomes parallel with cycle lengths $(b,a)$ and conversely.  In both
cases $\nu'=\nu$, the candidate pairs are the same for $t\le t_0$, and $F$ toggles.  The
column pairs with $t<t_0$ are disjoint from $\{q,c\}$, so their $\tau_t$, hence their
badness, are unchanged.  At $t_0$ the turn replaces $\tau_{t_0}$ by $\tau_{t_0}$ with the
cycle $Z$ reversed (or by the inverse of that, if the ordered pair is reversed), so $x_1,y_1$
are still in different cycles, $Z$ is still a cycle with the same intersection with
$\{1,2\}$, and the other cycle through $\{x_1,y_1\}$ is unchanged: $t_0$ and $Z$ are the
same in the turned square, and the turn is an involution.  The identity follows, and
$\Pr[F]\ge\tfrac12\Pr[G]$, $\Pr[F^c]\ge\tfrac12\Pr[G]$ give the bound.
\end{proof}

\begin{proposition}[the first ladder pair controls the bias]\label{prop:firstpairbias}
For every $n,\lambda,\alpha,\beta$,
\[
\Pr_X[B]-\Pr_X[A]=\E_X\bigl[(\mathbf 1_B-\mathbf 1_A)(\mathbf 1_{F^c}-\mathbf 1_F);\,G^c\bigr];
\]
consequently $\bigl|\Pr_X[B]-\tfrac12\bigr|\le\tfrac12\Pr_X[G^c]$,
$\Pr_X[A],\Pr_X[B]\ge\tfrac12\Pr_X[G]$, and
$|\bar q(\mu\to\lambda)-1|\le2\Pr_X[G^c]/(1-\Pr_X[G^c])$.
\end{proposition}

\begin{proof}
Write $B_{FG}=\Pr_X[B\cap F\cap G]$ and so on.  The trade $T_1$ of
Theorem~\ref{thm:ladder} ($k=1$) is an involution of $X\cap F$ whose frame is
$(x_1\,y_1)\circ\pi$, so it exchanges $A$ and $B$:
$B_{FG}+B_{FG^c}=A_{FG}+A_{FG^c}$.  The turn of Proposition~\ref{prop:firstpair} is an
involution of $X\cap G$ exchanging $F$ and $F^c$ and preserving $\pi$, hence $A$ and $B$,
which are functions of $\pi$: $B_{FG}=B_{F^cG}$ and $A_{FG}=A_{F^cG}$.  Hence
\begin{align*}
\Pr[B]-\Pr[A]&=(B_{F^cG}+B_{F^cG^c})-(A_{F^cG}+A_{F^cG^c})
=(B_{FG}-A_{FG})+(B_{F^cG^c}-A_{F^cG^c})\\
&=(A_{FG^c}-B_{FG^c})+(B_{F^cG^c}-A_{F^cG^c}),
\end{align*}
which is the displayed expectation; its absolute value is at most $\Pr[G^c]$.  ($T_1$ need
not preserve $G$.)  The last inequality is $|\bar q-1|=|\Pr[B]-\Pr[A]|/\Pr[A]$ with
$\Pr[A]\ge\tfrac12(1-\Pr[G^c])$.
\end{proof}

\begin{hypothesis}[first-pair hypothesis]\label{hyp:firstpair}
$\displaystyle\sup_{\lambda,\alpha,\beta}\Pr_{X(n,\lambda;\alpha,\beta)}[G^c]=o(1/\log n)$.
\end{hypothesis}

\begin{corollary}\label{cor:firstpair}
Hypothesis~\ref{hyp:firstpair} implies Conjecture~\ref{conj:cgw}, with
$\dtv(\PP_n,\QQ_n)=O(\log n\cdot\sup\Pr_X[G^c]+n^{-1+o(1)})$.
\end{corollary}
\begin{proof}
With $s=\sup\Pr_X[G^c]$, Proposition~\ref{prop:firstpairbias} gives
$|\bar q(\mu\to\lambda)-1|\le2s/(1-s)$ for every adjacent pair, i.e.\ $\mathrm{EQ}(\delta)$
with $\delta=2s/(1-s)=o(1/\log n)$, and Theorem~\ref{thm:reduction} applies.
\end{proof}

\paragraph{What the first-pair hypothesis asks, and what it costs.}  The event $G^c$ is
selected by the first ladder pair but is not local to it: it involves the column cycles of
the pairs $\{q_t,c_t\}$ through all rows, and rows $1,2$ through admissibility.  It is
again a lower tail of a count of events at square-selected positions --- the obstacle is
unchanged in kind --- and it decomposes as
\[
\Pr_X[G^c]\le\Pr_X[\nu\le M]+\Pr_X[\nu>M,\ t=1,\dots,M\text{ all bad}].
\]
The first term is an upper bound on a short arc or cycle of $\rho_{x_1,y_1}$ at $p,p'$;
nothing of the kind is proved at a frame-selected pair (Proposition~\ref{prop:splice}
bounds a different frame event), and heuristically it is $\approx4M/n$.  The second term
--- among $M$ candidates on disjoint free column pairs not all are bad --- is the
obstacle itself.  Compared with Hypothesis~\ref{hyp:witness}, at the level of heuristics
and data only: one row pair instead of $\ge n/\log^2n$; per-event probability $\approx0.4$
(the random-permutation value for ``$x_1\not\sim y_1$ under $\tau_t$ and one of the two
cycles has even intersection with $\{1,2\}$'') instead of $\ell'/n$; expected count
$\approx0.4\,\E\nu$ with $\E\nu=6.3,9.6,18$ at $n=30,50,100$ in the data
(Remark~\ref{rem:fixrect}), instead of $\ge\log^2n$ ($\log^3n$ at the parameters of
Proposition~\ref{prop:Lwitness}); no offset families, no splice bound, no conditioning on
$K_0$.  Neither hypothesis is known to imply the other: $\Pr_X[F^c]\ge\tfrac12\Pr_X[G]$, so
the first pair alone cannot give (L), and Hypothesis~\ref{hyp:witness} says nothing about
$G$.  So Hypothesis~\ref{hyp:firstpair} is a third alternative sufficient condition, the one
with the fewest moving parts; the reference hypothesis of the paper remains
Hypothesis~\ref{hyp:witness}.  The minimal instance \eqref{eq:minimal} of
\S\ref{sec:remains} follows from Hypothesis~\ref{hyp:firstpair} (with $c=\tfrac12-o(1)$)
but not evidently from Hypothesis~\ref{hyp:witness}, and is sufficient for neither.

\begin{remark}[data: completions of a fixed rectangle]\label{rem:fixrect}
To test the hypotheses on types that uniform sampling never produces we ran the
Jacobson--Matthews chain restricted to the completions of a fixed $2\times n$ rectangle
(\texttt{jmfix.c}: moves touching rows $1,2$ are never proposed from a proper state and are
rejected from an improper one; the walk is a simple random walk on the restricted move
graph, every proper state having restricted degree $(n-2)n(n-3)$, so its stationary law is
uniform on the proper states of the component of the start).  Connectivity of the restricted
graph is not proved; at $n=5,6$ (types $(5)$, $(6)$, $(2,2,2)$; $36$, $4032$, $5376$
completions, enumerated exactly) every completion is reached from two different starts with
$\chi^2$ statistics within one standard deviation of uniformity, and at $n=7$ the sampled
$\Pr[B\mid R]=0.4965$ agrees with the exact $0.497076$ over all $6566400$ completions of
$R=(\mathrm{id},(1\,2\cdots7))$.  Mixing is not analysed; runs with different thinning
and starts agree within noise.  Any two rectangles of type $\lambda$ are isotopic by a
column permutation and a symbol relabelling fixing rows $1,2$, which preserve completion
counts and carry a mark of class $(m,\alpha)$ to one of the same class; every event
considered here is defined through the cycle structures of row and column permutations and
is invariant under such isotopies.  Hence $\Pr_X$ of such an event equals its probability
under a uniform completion of any one rectangle $R$ of type $\lambda$ with any one mark of
the class, and averaging over the marks of the class in $R$ (which are equivalent under the
isotopies preserving $R$) only reduces the variance.  Findings (assuming connectivity and
mixing of the sampler) ($n=30$: $6\cdot10^4$ squares per type; $n=50$: $4\cdot10^4$;
$n=100$: $1.2\cdot10^4$; types $(n)$, $(n/2,n/2)$, $(n-2,2)$, $(10,10,10)$, $(4,2^{13})$,
$(6,3^8)$, $(4,2^{23})$, $(6,3^{14},2)$, $(4,2^{48})$; marks at $\alpha=2$,
$\approx m/3$, $\lfloor m/2\rfloor$):
(i) $\Pr_X[B]\in[0.4974,0.5024]$ in every class at $n\le50$ (sampling error
$\approx0.0025$), i.e.\ by Theorem~\ref{thm:marked}
$\CC_n(\lambda)/\CC_n(\mu)\in[0.995,1.005]$ for every tested adjacent pair, including
$(4,2^{13})\to(2^{15})$ and $(6,3^8)\to(3^{10})$, types of $\PP_n$-probability
$e^{-\Theta(n)}$; $[0.495,0.504]$ at $n=100$ (error $\approx0.005$).
(ii) $\Pr_X[F]\in[0.4988,0.5017]$ in every class at $n\le50$ ($0.4896$ at $n=7$, so
$\tfrac12$ is not an identity).
(iii) Along the ladder, $\Pr[\text{no witness among the first }K\text{ pairs}\mid K_0-1\ge K]$
equals $(1-w_1)^K$ within sampling error for $K\le12,20,40$ at $n=30,50,100$ and
thresholds $\ell'=4,8,16,24$, where $w_1$ is the first-pair witness rate ($0.100$, $0.060$,
$0.030$ at $\ell'=4$); the hazard $\Pr[\text{witness at }k\mid\text{none before }k]$ is
flat in $k$ to three digits; $\Pr[\text{none of the first }K\text{ ladder pairs flippable}\mid K_0-1\ge K]=2^{-K}$
within noise to $K=12$--$15$ ($7\cdot10^5$ instances at $n=30$); across the diagonals $D_c$ of
$B$-instances with $d_{12}=\ell\le n/4$, which share their $y$-rows, the number of
diagonals carrying a witness vanishes with the binomial probability and has binomial
lower tails, for all $\ell\le12$.
All of this is identical for the extreme types.
(iv) For Hypothesis~\ref{hyp:firstpair}, with all candidates $t<\nu$: $\Pr_X[G^c]=0.61,
0.18, 0.11, 0.055$ at $n=7,30,50,100$ (type $(n)$, $\alpha=2$; $0.056$ for $\alpha=50$ at
$n=100$), i.e.\ $\approx5.5/n$ for $n\ge30$; $\E\nu=6.3,9.6,18$ and $\Pr[\nu=k]\approx4/n$
for small $k$; $\Pr[G^c\mid\nu=k]=0.59,0.37,0.22,0.13,0.08$ for $k=2,\dots,6$ at $n=30$
against $0.59^{k-1}$ (a slight positive dependence among the candidates, which share the
rows $x_1,y_1$); $t=1$ is good with probability $0.41$ in every class at $n\le100$.  Other
types were measured only with $t\le8$: $\Pr[\text{some }t\le8\text{ good}]=0.81,0.88,0.94$
at $n=30,50,100$ in every class.  The identity of
Proposition~\ref{prop:firstpairbias} and its two ingredients hold to all digits on the
exact enumerations at $n=5,6$ (every class) and within sampling error at $n=7$ ($-0.0052$
and $-0.0058$, each $\pm0.001$, against the exact $-0.005848$).
None of these data reach the regime of any hypothesis ($\ell'\ge\log^5n$, failure
$o(1/\log n)$); they exclude a visible dependence between the events, or a visible effect
of the type, at the first- and second-moment level for $n\le100$.
\end{remark}

\section{What remains}\label{sec:remains}

Hypotheses~\ref{hyp:witness}, \ref{hyp:orbitgen}, \ref{hyp:adaptive}
and~\ref{hyp:firstpair} are statements about configurations at positions selected by the
frame: a short $\rho_{x_k,y_k}$-cycle through $p$ avoiding $p'$; a clean column pair
separated by the mark; a short clean column cycle at a prescribed column pair; a column pair
along $\rho_{x_1,y_1}$ whose column permutation separates $x_1$ from $y_1$.  Their heuristic
frequencies are, for large $n$, far above what is needed; the difficulty is entirely that
the positions are selected by the frame of the mark and the whole space is conditioned on
the type of rows $1,2$.  The minimal instance of the obstacle is a lower bound on a single
frame-selected pair:
\begin{equation}\label{eq:minimal}
\Pr_X\bigl[(x_1,y_1)\text{ flippable}\bigr]\ge c,\qquad x_1=\pi^{-1}(1),\ y_1=\pi^{-1}(2),
\end{equation}
which no identity of this paper gives unconditionally: flippability of a ladder pair is a
class invariant of the ladder trades, and Proposition~\ref{prop:firstpair} reduces
\eqref{eq:minimal} one way to $\Pr_X[G]\ge2c$, a lower bound for an event of the same kind
with rows and columns exchanged (the column permutation of a pair of columns selected from
rows $x_1,y_1$ separating the frame-selected rows $x_1,y_1$), which is not known to be
easier.  By row exchangeability \eqref{eq:minimal} is a statement at fixed positions under
conditioning: it is the probability that $p,p'$ lie in different cycles of $\rho_{3,4}$
given rows $1,2$ and the two cells $L(3,p)=L(1,p')$, $L(4,p)=L(2,p')$.  Even without the
conditioning on rows $1,2$ we know of no lower bound in the literature for the related
quantity $\Pr[\text{the }\rho_{3,4}\text{-cycle through }p\text{ has length }k\text{ and avoids }p']$,
$3\le k\le\log^5n$; the one case we know to be settled is $k=2$, where the concentration
of the number of intercalates at $n^2/4$ \cite{KSS21} and the exchangeability of row
pairs and of columns give $\approx1/n$.  We record what
the available tools do and do not give.

\emph{Unconditional facts.}  For two generic rows of a uniform square,
Proposition~\ref{prop:tailsurvive}(i) gives, by Markov's inequality on the number of
columns on short cycles, $\Pr[\text{all cycles of }\rho_{x,y}\text{ have length}\le\varepsilon n]\le2\varepsilon/(1-\varepsilon)$;
so a generic pair has a cycle longer than $n/5$ with probability $\ge\tfrac12$, and the
column-cycle analogue holds by transposition.  Kwan and Sudakov's bound
$\Pr[N_2\ge t]=e^{-\Omega(t\log t)}$ on the number of $2$-cycles of two fixed rows
\cite[Thm.~3]{KS18} is proved by joining switchings of the same kind as those of
Proposition~\ref{prop:tailsurvive} and is likewise an upper bound on short cycles.  The
number of completions of any $2\times n$ rectangle is within $e^{O(n\log^2n)}$ of any other
(van der Waerden and Br\'egman; \cite[Prop.~5]{KS18}), so any property of uniform squares
with failure probability $e^{-\omega(n\log^2n)}$ --- such as the intercalate concentration
of Kwan, Sah and Sawhney \cite{KSS21} --- holds conditionally on any fixed rows $1,2$ and
mark, hence in $X$ for every $\lambda,\alpha,\beta$; but the natural failure probability of a
statement about $m\le n$ prescribed row pairs is $e^{-\Omega(m)}$, which this transfer does
not reach.

\emph{Why the fixed-matching version is harder.}  With a fixed matching the orbit weight
$\bar y_k$ of Proposition~\ref{prop:orbit} needs a lower bound on the probability that a
generic column pair is interleaved with $\{p,p'\}$ in $\rho_{x_k,y_k}$ --- an instance of
the statement ``two given pairs of lines of a random Latin square are interleaved in the
cycle structure of a third pair with probability in $[c,1-c]$'' inside $X$.  Without the
conditioning its upper half is elementary and its lower half reduces to macroscopic cycles
(by conjugation invariance the interleaving probability given the cycle type of $\rho_{x,y}$
is $\tfrac13\sum_i(m_i)_4/(n)_4+2\sum_{i\ne j}(m_i)_2(m_j)_2/(n)_4\le\tfrac13+O(1/n)$, and
$\ge\tfrac13\Pr[W\mid\text{type}]^3-O(1/n)$ by Jensen, $W$ the event that two given columns
share a cycle).  Inside $X$ the exact tools are the free trades of rows other than $1,2$
(which preserve interleaving) and the legal column turns (which are legal exactly off the
interleaved configurations), and every identity they produce is fair: it equates two
expectations with random weights whose lower bounds are again of the same form.  Lower
bounds on ``different cycles'' events are, in the language of \cite{CGW08}, of the type of
the direction of their Lemma~3.12 whose proof has the gap: they require a repair of the
same-cycle configurations with a rigid inverse, and the repairs available in $X$ (trades or
cross-switches of a free pair of third rows) have unbounded multiplicity.  The one case in
which we found a rigid inverse is CGW's own case 3 with $\min(\alpha,\beta)=2$, which gives
$\Pr_X[A]\ge\tfrac14$ for marks at distance $2$ (Remark~\ref{rem:gap}); it does not extend
to other distances for a reason of type rather than multiplicity.  The adaptive
coordinates sidestep the question: the coordinate is placed where it is separated by
construction, and the price is paid in rarity rather than in a constant we cannot prove;
at the first ladder pair alone (\S\ref{sec:orbit}(e)) the rarity disappears as well, and
what is left is an upper bound on a short arc at the frame-selected pair together with the
lower tail of the number of good candidates among the $\nu-1$ column pairs along
$\rho_{x_1,y_1}$ --- the smallest sufficient condition we have, and still a lower bound at
square-selected positions.


\begin{thebibliography}{9}
\bibitem{CGW08} N.~J. Cavenagh, C. Greenhill and I.~M. Wanless, The cycle structure of two
rows in a random Latin square, \emph{Random Structures Algorithms} 33 (2008), 286--309.
\bibitem{gapnote} [author], A gap in the proof of Lemma 3.12 of Cavenagh--Greenhill--Wanless, note, October 2026.
\bibitem{JM96} M.~T. Jacobson and P. Matthews, Generating uniformly distributed random Latin
squares, \emph{J. Combin. Des.} 4 (1996), 405--437.
\bibitem{KS18} M. Kwan and B. Sudakov, Intercalates and discrepancy in random Latin squares,
\emph{Random Structures Algorithms} 52 (2018), 181--196.
\bibitem{KSS21} M.~Kwan, A.~Sah and M.~Sawhney, Large deviations in random Latin squares,
\emph{Bull. London Math. Soc.} 54 (2022), 1420--1438.
\bibitem{repo} [author], \texttt{cgw-conjecture}: scripts, logs and notes, \url{https://github.com/mkkinyon/cgw-conjecture}, 2026.
\end{thebibliography}
\end{document}
